\documentclass[a4paper,12pt]{article}
\usepackage[UTF8]{ctex}
\usepackage{fullpage}
\usepackage{iftex}
\ifPDFTeX
  \usepackage[T1]{fontenc}
  \usepackage{mathptmx}
  \usepackage[scaled]{helvet}
  \usepackage{courier}
\else
  \usepackage{fontspec}
  \usepackage{unicode-math}
  \setmainfont{Times New Roman}
  \setsansfont{Arial}
  \setmonofont{Consolas}
  \setmathfont{Cambria Math}
\fi
\usepackage{textcomp}
\usepackage{graphicx}
\usepackage{amsmath}
\usepackage{cite}
\usepackage{url}

\begin{document}

\title{基于 nnU-Net 的视网膜血管分割方法研究}
\author{}
\date{\today}

\maketitle

\begin{abstract}
本文围绕基于 nnU-Net 的视网膜血管分割方法展开研究，针对细小血管恢复不足、跨层信息利用有限以及结构连续性表达不稳等问题，在统一实验协议下对模型结构与结果优化策略进行了逐步探索。研究首先通过前期实验考察输入增强与损失函数调整的作用，结果表明，在 nnU-Net 既有预处理与训练流程基础上，额外输入端预处理和损失函数调整都没有带来明确且稳定的最终收益，后续研究因此更多转向网络结构本身。在后期主实验中，本文以 nnU-Net 为统一基线，围绕 PA-Net 模块迁移、LPT 主线筛选以及可变形采样、多尺度建模、位置注意力和连续几何轨迹建模依次推进，最终形成了以 MS-DSConv + CA 为核心并辅以轻量后处理的分割方案。在后期五折验证协议下，最终方案的 F1-score 达到 0.8027，较后期 nnU-Net 基线的 0.7891 小幅提升 0.0136。结果提示，在本文采用的统一协议下，围绕细长结构连续几何轨迹展开的结构改进与细小血管恢复和结构表达的小幅稳定改善相关，同时能够维持检出能力与误检控制之间的可接受平衡。
\end{abstract}

\section*{致谢}
\addcontentsline{toc}{section}{致谢}
首先，我要衷心感谢电子科技大学数学科学学院于佳丽教授在本次毕业设计过程中给予我的指导与帮助。在课题开始阶段，我对医学图像分割的了解还比较有限，也不太清楚一个完整的研究课题应当如何规划和推进。于老师帮助我确定了合适的选题，明确了研究方向，并耐心地引导我逐步进入这一领域。在后续写作和修改过程中，于老师也提醒我不仅要关注实验结果本身，还要重视实验协议的可靠性、文章逻辑的完整性以及结论表述的分寸感。

同时，我要特别感谢课题组博士生刘志帅师兄在整个研究过程中的耐心帮助。无论是 nnU-Net 框架的使用、消融实验的设计，还是对不稳定实验结果的分析，对我来说一开始都并不轻松。每当实验结果没有达到预期，或者某个改动方向暂时走不通时，刘师兄都会和我一起分析可能原因，并提出更实际的调整思路。这些讨论让我逐渐意识到，毕业设计并不是简单地追求更高的指标，而是要判断一个改动是否稳定、是否有解释性，以及是否真正服务于研究问题。

此外，感谢课题组其他博士生同学以及相关老师、技术人员在讨论、技术细节和实验环境方面给予的帮助。在接触不熟悉的工具、训练设置和评价流程时，他们的建议减少了我很多困惑。课题组中较好的交流氛围也让我看到，研究问题往往需要在不断讨论和修改中逐步变得清晰。

我也感谢电子科技大学以及格拉斯哥学院各位老师在本科阶段给予我的培养。过去几年的课程学习为我完成本课题提供了必要的数学基础和计算能力，而毕业设计则让我有机会把课堂中学到的知识与一个具体的应用研究问题联系起来。通过这个过程，我也更加清楚地认识到，方法实现、实验验证和论文表达之间并不是彼此独立的环节，而是共同构成了一项完整工作的不同部分。

最后，感谢我的家人、同学和朋友在这段时间里给予我的理解、鼓励与支持。在持续实验和反复修改论文的过程中，他们的陪伴让我能够保持相对稳定的心态。本次毕业设计对我来说既有挑战，也很有意义。感谢所有在这个过程中给予我帮助的人，这些经历也会继续影响我之后的学习和工作。

\tableofcontents
\clearpage

\section{绪论}

\subsection{研究背景与意义}
视网膜是人体中少数可以通过无创成像直接观察微循环状态的组织。正因如此，眼底图像中的血管信息长期被用于疾病筛查、辅助诊断和病情评估 \cite{staal2004drive,hoover2000locating}。在临床实践中，糖尿病视网膜病变、高血压、动脉硬化以及部分心脑血管相关疾病，往往都会伴随视网膜血管口径、分支形态、弯曲程度和连续性的变化。若能较准确地完成血管分割，就可以为后续的血管宽度测量、动静脉分析、病灶定位和定量评估提供基础，也有助于提高眼科辅助诊断系统的自动化水平和稳定性 \cite{fraz2012ensemble,budai2013robust}。

不过，这个任务本身并不轻松。视网膜血管天然具有细长、弯曲和多尺度分布的特点，尤其是末梢微细血管，宽度小、对比度低，还容易受到成像噪声、照明不均和背景纹理的干扰 \cite{staal2004drive,hoover2000locating}。与此同时，血管分叉、交叉以及局部断裂区域又会进一步增加建模难度。如果分割结果不能较好地保持血管连续性和拓扑一致性，后续定量分析的可靠性就会明显下降。因此，如何在保证整体分割准确率的同时，更有效地恢复细小血管并维持结构连续性，一直是该任务中的核心问题。

近年来，深度学习方法已经成为医学图像分割研究的重要方向。以 U-Net 为代表的编码器—解码器结构，凭借端到端训练方式和较强的像素级预测能力，在多种医学图像分割任务中表现出较好的适应性 \cite{ronneberger2015unet,cicek2016unet3d,milletari2016vnet}。在此基础上，nnU-Net 进一步将预处理、网络配置、训练策略和后处理步骤进行了系统化的自适应设计，降低了模型构建和参数调节的门槛，因此逐渐成为医学图像分割研究中常用的通用基线框架 \cite{isensee2021nnunet}。对于视网膜血管分割这样既要求全局语义一致性、又要求细粒度结构恢复的任务来说，nnU-Net 提供了一个比较稳定且便于扩展的起点。

但通用框架并不意味着已经足够适配具体任务。视网膜血管属于典型的细长目标，模型最终能否分得好，不只取决于多尺度特征提取，还取决于它能否更好地处理长程依赖、关键空间位置、跨层信息交互以及几何连续性等问题。已有研究表明，注意力机制、特征融合策略以及更灵活的采样方式，确实有助于改善传统编码器—解码器结构在细小血管恢复、边界保持和结构连续性方面的不足 \cite{luo2025panet,oktay2018attentionunet,zhu2021deformabledetr,hou2021coordatt}。因此，在 nnU-Net 框架基础上进一步探索面向视网膜血管特征的结构改进，既有实际应用意义，也有一定的方法研究价值。

基于以上考虑，本文以 nnU-Net 为基础框架，围绕视网膜血管分割中细小血管检测不足、跨层信息利用有限以及结构连续性建模不充分等问题展开研究。通过对不同结构模块、采样机制、注意力策略和输出优化方法的逐步探索，本文希望为构建更加适合视网膜血管分割任务的模型提供实验依据和方法参考。

\subsection{国内外研究现状}
视网膜血管分割是眼底图像分析中的基础问题之一。相关研究最早主要从图像增强和规则提取入手，例如匹配滤波、阈值分割、区域生长和形态学处理等方法。这类方法在血管响应较清晰、成像条件相对稳定时能够取得一定效果，但对图像质量、噪声和参数设置往往比较敏感，一旦遇到复杂背景、低对比度区域或细小血管，性能就容易下降 \cite{hoover2000locating,staal2004drive}。后来，随着公开数据集和评价标准逐步完善，研究者开始引入机器学习和集成分类方法来提高血管与背景之间的区分能力，不过这一路线在很大程度上仍然依赖人工设计特征 \cite{fraz2012ensemble,budai2013robust}。

深度学习的出现明显改变了这一局面。卷积神经网络能够直接从原始图像中学习层次化特征，在医学图像分割任务中表现出更强的适应性。尤其是 U-Net 及其编码器—解码器结构，通过下采样提取高层语义信息，再借助跳跃连接恢复空间细节，在多种医学图像分割任务中都取得了较好效果 \cite{ronneberger2015unet,cicek2016unet3d,milletari2016vnet}。在视网膜血管分割领域，这类方法也很快取代了早期传统方法，成为后续研究的主要基础 \cite{liskowski2016segmenting,maninis2016driu,oliveira2018fcn}。

不过，标准 CNN 和 U-Net 并没有把这个问题彻底解决。视网膜血管并不是规则、厚实的目标，而是细长、弯曲、跨尺度延伸的结构。只依赖局部卷积时，模型往往能较好识别主干血管，但在细小血管、弱响应区域和复杂分叉附近仍然容易出现漏检或断裂。另一方面，标准编码器—解码器结构中的信息交互主要依赖固定跳跃连接，对于远距离依赖、关键空间位置以及跨层结构连续性的刻画仍然有限。也正因为如此，即便 U-Net 已经取得了较好效果，围绕其局限继续做针对性改进仍然很有必要。

后续不少工作并没有脱离 U-Net 的基本框架，而是在其基础上做更细的调整。有的工作重点放在编码器与解码器之间的连接方式和多尺度融合上，例如 UNet++ 通过嵌套式跳跃连接缩小两者之间的语义差距 \cite{zhou2018unetpp}。也有工作选择在网络中引入注意力机制，例如 Attention U-Net 利用注意力门抑制无关响应、突出目标区域 \cite{oktay2018attentionunet}。这些研究说明，U-Net 作为基础框架依然有效，但在具体任务中，特征融合方式、空间选择能力和结构恢复能力往往都需要进一步调整。

对视网膜血管分割来说，这一点尤其明显。血管是连续而细长的结构，方向性也很强，仅靠局部卷积并不容易利用那些相距较远、但在拓扑上彼此相关的特征。弱响应区域往往正好出现在这些地方，模型如果处理不好，就很容易把一条本来连通的血管分成几段。正因如此，注意力机制、多尺度融合以及更灵活的采样方式，逐渐成为这一方向里比较常见的改进思路。

nnU-Net 的出现又让医学图像分割研究进一步走向标准化。与其说 nnU-Net 是一套单独的新网络，不如说它是一套面向任务的自适应配置方案，能够根据具体数据集自动调整预处理、网络拓扑、训练计划和后处理流程 \cite{isensee2021nnunet}。由于它在多个医学图像分割任务中都表现出较强鲁棒性，也不容易被轻易超越，因此已经成为这一领域常用且稳定的基线框架。对本文而言，选择在 nnU-Net 上开展改进，正是因为它的起点足够高，后续结构调整是否有效也就更容易被检验出来。

近年的视网膜血管分割研究，开始更直接地围绕特征融合和采样机制展开。PA-Net 是其中较有代表性的一项工作。该方法在 U-Net 风格主干上引入 Lightweight Parallel Transformer (LPT) 和 Adaptive Vascular Feature Fusion Module (AFFM)，并配合针对眼底图像的预处理流程，试图同时加强全局结构建模和多层级血管特征融合 \cite{luo2025panet}。这类工作也提醒我们，仅靠标准编码器—解码器结构往往不足以稳定恢复细小血管，模型还需要在关键位置建模和结构连续性表达上更进一步。

除了网络结构本身，近年来也有不少工作开始关注细长结构的拓扑一致性。深监督可以借助多层辅助输出强化多尺度特征学习 \cite{lee2015dsn}；可变形注意力和调制机制通过学习采样偏移与权重，使模型在空间上具有更大的灵活性 \cite{zhu2021deformabledetr,zhu2019dcnv2}；Coordinate Attention 通过方向感知的方式补充位置信息表达 \cite{hou2021coordatt}；Dynamic Snake Convolution 则尝试用连续轨迹式采样去贴近细长目标的真实几何形态 \cite{qi2023dsc}。与此同时，围绕血管连通性设计的损失函数和输出后处理策略，也在被用于弥补结构建模上的不足。研究重点已经不再只是“能不能把血管分出来”，而是进一步转向“如何更稳定地保留细小血管、维持结构连续性，并尽量控制误检”。

现有研究已经为视网膜血管分割提供了较为丰富的方法基础，但在细小血管检出、跨层信息利用、拓扑连续性保持以及综合指标平衡等方面，仍然有继续改进的空间。基于这一点，本文选择以 nnU-Net 作为统一基础框架，在已有工作的启发下，进一步探索更适合视网膜血管分割任务的模型设计与优化方法。

\subsection{视网膜血管分割面临的主要问题}
视网膜血管分割看起来像一个像素级分类任务，真正困难的地方却远不止于此。与一般器官或病灶分割相比，它更依赖细粒度结构恢复和整体拓扑保持，因此问题通常不会单独出现，而是同时落在成像条件、特征表达和结果判别几个层面。

最直接的困难来自血管本身。血管是典型的细长结构，弯曲明显，而且跨尺度分布。主干血管通常较容易识别，真正难的是末梢微细血管、弱响应区域和复杂分叉位置。这些区域宽度小、对比度低，也更容易与背景纹理混在一起 \cite{staal2004drive,hoover2000locating}。一旦模型在这里处理得不够细，结果往往不是简单少分一点，而是出现局部断裂、细小分支缺失，甚至把原本连续的血管切成几段。

成像条件和数据差异又把这个问题进一步放大。不同设备、不同采集环境以及不同个体之间，眼底图像在亮度、对比度、噪声水平和背景纹理上往往都有明显差别 \cite{fraz2012ensemble,budai2013robust}。公开常用的眼底数据集数量也相对有限，这使得模型泛化能力的验证更受数据条件约束。即便在同一数据集中，视盘高亮区域、病灶反光、局部阴影和颜色分布变化，也可能干扰血管响应。模型面对的并不是规则而干净的目标，而是在复杂背景中识别细长结构，这对鲁棒性提出了更高要求。

另一个难点出现在特征建模层面。卷积神经网络擅长提取局部纹理信息，但对远距离依赖关系、关键空间位置以及跨层结构连续性的刻画仍然有限。标准编码器—解码器结构虽然能够借助跳跃连接补回一部分空间细节，但在处理细长、跨尺度延伸的血管时，仍然容易出现高层语义足够、浅层细节不够，或者局部响应存在但整体连通性不佳的情况。围绕这一问题，如何更有效地组织跨层信息交互，并提高模型对关键位置和连续结构的感知能力，已经成为这一领域持续关注的方向 \cite{zhu2021deformabledetr,zhu2019dcnv2,hou2021coordatt}。

评价模型时还会遇到另一组经常并存的矛盾，即检出能力与误检控制之间的平衡。对细小血管和弱响应区域更敏感，通常有助于提高灵敏度，但同时也可能把部分噪声或背景纹理误判为血管，从而拉低精确率。反过来，如果模型过于追求结果干净，又容易牺牲细小血管的恢复能力。就这一任务而言，单一指标往往不足以说明模型是否真的更好，通常还需要结合灵敏度、精确率、特异性和 F1-score 一起判断。

在此基础上，另一个现实问题是：即使已经有 nnU-Net 这样的成熟基线，稳定取得改进也并不容易。nnU-Net 已经提供了较为成熟的预处理、训练与推理流程 \cite{isensee2021nnunet}，后续改进如果要体现价值，往往不能只停留在表面的模块叠加或局部参数调整上，而需要更有针对性地回应前面这些问题。近年的一些工作也正是在细长结构的空间建模、跨层信息利用和输出结果约束等方面继续推进 \cite{luo2025panet,qi2023dsc,zhu2021deformabledetr,hou2021coordatt}。

总的来说，视网膜血管分割的几个核心问题并不陌生：细小血管和弱响应区域难以稳定恢复，复杂成像条件下的鲁棒性不足，局部细节与全局结构难以同时兼顾，以及检出能力和误检控制之间始终存在权衡。本文后续讨论模型设计与优化策略时，也将主要围绕这些问题展开。

\subsection{本文的研究内容与技术路线}
围绕前文所述的几个主要问题，本文以 nnU-Net 为统一基础框架，针对视网膜血管分割中细小血管恢复不足、跨层信息利用不充分以及结构连续性保持困难等现象，开展模型改进与结果优化研究。研究并不试图脱离现有成熟框架重新搭建模型，而是在统一实验条件下考察不同改动对分割性能的实际影响，据此筛选更适合当前任务的结构设计与优化策略。

在研究推进过程中，本文首先建立基线模型，并结合输入增强与损失函数调整进行前期探索，以判断在 nnU-Net 既有预处理与训练流程基础上，这类额外改动究竟还能带来多大增益。在此基础上，后续实验统一在同一数据集和一致的训练协议下展开，使不同模型之间的比较尽可能建立在相同条件之上，从而保证后续分析的可解释性 \cite{isensee2021nnunet}。

从具体内容来看，本文的工作不是同时铺开，而是沿着几条相互衔接的线索逐步推进。前一部分主要建立统一参照并完成前期可行性探索，包括基线构建、输入增强测试和损失函数调整，用来判断在 nnU-Net 已有内部预处理流程基础上，额外输入端处理和相关训练设置调整是否还能提供稳定收益。随后，研究重心转向网络结构本身，通过迁移、比较和筛选相关模块，考察它们在当前框架下的适配性，并据此决定哪些思路值得继续保留。待结构主线逐渐清晰之后，再从输出层面对结果进行补充优化与约束，使模型不仅在检测能力上有所提升，也在整体平衡性和结果稳定性上表现更稳。

在结构改进方面，本文先参考已有工作的设计思路，对相关模块进行迁移与适配，并通过消融比较它们在当前框架中的实际收益。随后，研究重点逐步转向高层特征建模与跨层信息交互，围绕更灵活的空间采样、更有效的多尺度特征利用以及更适合细长结构表达的建模方式持续推进。这条技术路线并不是简单地叠加模块，而是在不断比较和调整中逐步收拢到相对稳定的改进方向 \cite{luo2025panet,zhu2021deformabledetr,zhu2019dcnv2,hou2021coordatt,qi2023dsc}。

除结构设计外，本文还从输出结果优化和损失函数约束两个方面进行了补充探索，包括后处理策略以及面向拓扑一致性的附加尝试。这样安排的目的，不是单独追求某一项指标更高，而是希望在细小血管恢复、结构连续性保持与误检控制之间找到更合适的平衡。整体上，本文形成了一条由基线建立、结构改进到结果优化逐步推进的技术路线，后续章节也将围绕这条路线依次展开理论基础、模型实现和实验分析。

\subsection{论文结构安排}
全文共分为五章。第一章为绪论，交代视网膜血管分割的研究背景与意义，梳理国内外研究现状，说明该任务面临的主要问题，并给出本文的研究内容与技术路线。第二章为理论基础与相关方法，介绍视网膜血管分割任务、U-Net 与 nnU-Net 的基本原理，以及后续模型改进涉及的相关方法。第三章为改进模型设计与实现，沿着本文的实验主线展开，说明基线建立、结构迁移与筛选、主线结构改进和结果优化过程。第四章为实验结果与分析，对各阶段模型的实验结果进行集中比较，并结合关键指标变化讨论模型性能和结构演化路径。第五章为总结与展望，对全文工作加以归纳，同时分析当前研究的不足与后续可能的改进方向。

\section{理论基础与相关方法}

\subsection{视网膜血管分割任务概述}
视网膜血管分割的任务，是从眼底图像中识别并提取血管区域，得到与原图对应的血管掩膜。形式上，它通常被归为像素级分类问题，即判断每一个像素属于血管还是背景；但在实际应用中，这项工作并不只是一次前景提取。分割结果往往还要用于后续的血管宽度测量、动静脉分析、病灶定位以及眼底定量评估，因此除了像素判别精度之外，血管结构是否连续、分支是否完整，同样会直接影响结果的使用价值 \cite{staal2004drive,hoover2000locating}。

现有研究通常以彩色眼底图像作为输入，输出与原图空间位置对应的血管掩膜。掩膜一般由人工标注或专家修订得到，并作为训练和评价时的参考标准。与一般目标分割相比，视网膜血管在图像中所占面积通常较小，而背景区域却占据绝大多数像素，因此类别不平衡几乎是这一任务的固有问题。再加上血管本身细长、弯曲、分叉频繁，尺度变化也比较明显，模型既要识别局部纹理，又不能轻易破坏整体树状结构。

公开研究中常用的数据集主要包括 DRIVE、STARE、CHASE\_DB1 和 HRF 等 \cite{staal2004drive,hoover2000locating,chasedb12013,budai2013robust}。这些数据集为视网膜血管分割研究提供了较为常见的测试基础，但它们在成像设备、图像分辨率、视场范围、病理表现以及标注习惯等方面并不完全一致。因此，在阅读相关研究结果时，通常仍需结合具体数据条件理解不同方法的表现差异。

评价结果时，视网膜血管分割通常需要结合灵敏度、特异性、准确率、精确率以及 F1-score 等指标进行综合衡量 \cite{staal2004drive,hoover2000locating}。其中，灵敏度更能反映模型对细小血管和弱响应区域的检出能力，F1-score 更适合观察检出能力与误检控制之间的总体平衡。相比之下，精确率与特异性更多体现模型抑制伪阳性的能力。就本文关注的血管恢复和误检控制问题而言，后续分析会更重视灵敏度、精确率与 F1-score 的变化，同时结合特异性和准确率判断整体预测平衡。

归根结底，这一任务对模型提出的是几项并行存在的要求：既要尽量恢复细小血管和复杂分叉区域，减少漏检，又要避免把背景纹理或噪声误判为血管，同时还应尽可能保持血管网络在空间上的连续延伸。后续围绕网络结构、特征交互方式以及结果优化策略展开的研究，基本都可以看作是对这些任务要求的进一步回应。

\subsection{U-Net 与 nnU-Net 基本原理}
U-Net 是医学图像分割领域中最具代表性的编码器—解码器结构之一 \cite{ronneberger2015unet}。它提出时，卷积神经网络已经在图像分类任务中取得了较好效果，但医学图像分析更关注像素级定位，而不是单一类别标签。另一方面，医学图像任务通常难以获得大规模标注样本，直接照搬自然图像分类网络并不现实。U-Net 的价值，正在于它把上下文语义建模和精细空间定位放进同一套端到端框架中，使有限标注条件下的分割任务也能得到较稳定的结果。

从结构上看，U-Net 由收缩路径和扩张路径两部分组成。收缩路径通过多层卷积与下采样逐步扩大感受野，用来提取更高层次的语义信息；扩张路径则通过上采样逐步恢复空间分辨率，使特征重新回到与输入图像对应的尺度 \cite{ronneberger2015unet}。其中最关键的是跳跃连接，即把编码阶段较浅层的特征直接传递给解码阶段对应层级，用来补回上采样过程中容易丢失的空间细节。对医学图像分割来说，这一点很重要，因为目标边界、细小结构以及局部纹理往往主要保留在浅层特征中，而深层特征虽然语义更强，却更容易牺牲精细定位能力。U-Net 的整体结构示意如图 \ref{fig:unet_arch_cn} 所示。

\begin{figure}[htbp]
    \centering
    \includegraphics[width=0.82\textwidth]{pic/unet.png}
    \caption{U-Net 经典结构示意图。左侧为编码器路径，通过卷积和下采样逐步提取高层语义特征；右侧为解码器路径，通过上采样逐步恢复空间分辨率；中间的跨层连接用于将浅层细节信息传递到对应解码阶段。图片来源：文献 \cite{ronneberger2015unet}。}
    \label{fig:unet_arch_cn}
\end{figure}

也正因为编码器、解码器与跳跃连接之间形成了这种互补关系，U-Net 在医学图像分割任务中表现出较强的适应性。它既能利用下采样获得相对稳定的上下文信息，避免模型只依赖局部纹理作出判断，又能通过跨层特征融合尽量保留目标边界和细节。对于视网膜血管分割这样同时要求局部细节恢复与整体结构连续性的任务来说，U-Net 所体现出的这种设计思路至今仍具有基础性意义。

在 U-Net 之后，医学图像分割领域出现了大量结构变体，但这些方法往往需要针对具体数据集、硬件条件和训练策略做出较多人工调整。nnU-Net 的提出，正是针对这一现实问题 \cite{isensee2021nnunet}。与其说 nnU-Net 是一种全新的网络结构，不如说它是一套面向医学图像分割任务的自适应配置框架。它并不试图重新发明一个完全不同的主干网络，而是根据数据集特征自动决定预处理方式、网络配置、训练策略以及后处理方案，让分割流程尽可能标准化和可复用。

nnU-Net 的一个重要特点，是把任务相关信息尽可能前置为可分析的数据属性，再据此组织后续流程。原始论文中提出了类似“数据指纹”和“流程指纹”的思路，即先根据数据集的体素或像素间距、图像尺寸、类别分布以及目标区域大小等信息，对任务特征进行归纳，再进一步决定更合适的网络配置和训练方案 \cite{isensee2021nnunet}。在这种框架下，nnU-Net 不是先固定一套网络再强行适配所有任务，而是先理解数据集本身的形态特征，再反过来组织分割流程。

落实到具体流程中，nnU-Net 会自动组织多个关键环节，包括体素或像素间距处理、图像归一化、patch 尺寸设置、batch size 调整、网络深度与卷积核配置选择，以及 2D U-Net、3D full-resolution U-Net 和级联式 3D U-Net 等不同候选方案的生成与比较 \cite{isensee2021nnunet}。在训练阶段，它通常还会结合多层辅助输出的深度监督机制，以强化不同尺度特征的学习 \cite{isensee2021nnunet,lee2015dsn}。训练完成之后，它再结合验证结果对模型或模型集成方案进行筛选，并对推理后的输出执行必要的后处理，例如是否保留最大连通域等。由此可以看出，nnU-Net 的优势并不只体现在网络主体上，而是体现在“预处理—训练—验证—推理—后处理”这一整条分割流程的统一组织能力上。

这样一来，许多原本依赖经验的工程决策被转化为更系统的规则组合，研究者也能把更多注意力放在任务本身，而不是反复进行配置试错。也正因为 nnU-Net 已经在多类医学图像分割任务中表现出较好的稳定性，它不仅是一个常用基线，也可以被看作后续结构改进研究的可靠起点。对本文而言，这一点尤为重要，因为后续工作关注的并不是重新寻找一个新的起点，而是在一个统一且较稳定的框架之上，进一步考察结构改动是否能够真正带来增益。

对本文而言，U-Net 提供的是医学图像分割中的基本结构范式，而 nnU-Net 提供的则是在该范式基础上的统一实现与配置框架。前者说明了编码—解码结构和跳跃连接为什么能够成为医学图像分割的基本方案，后者则进一步说明，在已有成熟框架支持下开展改进研究，能够使后续比较建立在更稳定且更具可解释性的基础之上。基于这一认识，本文后续模型设计与实验分析均以 nnU-Net 为统一基础展开。

\subsection{注意力机制与特征融合方法}
在医学图像分割任务中，仅依赖卷积操作往往难以同时兼顾局部细节、长程依赖和跨层语义信息。对视网膜血管分割而言，这一问题更加明显：血管既有细小而密集的局部纹理，又具有连续延伸、分叉频繁和方向变化明显的整体结构。也正因为如此，模型既需要更准确地关注与目标相关的有效信息，也需要在不同层级之间更合理地组织特征表达。注意力机制与特征融合方法，正是在这样的背景下逐渐发展起来的。

注意力机制的基本出发点，是让模型在特征表示过程中对更重要的位置、通道或结构给予更高权重，从而抑制无关背景的干扰。相较于单纯依赖局部卷积的方式，注意力机制能够在一定程度上增强模型对关键区域的选择能力。以 Attention U-Net 为代表的方法，在 U-Net 结构中引入注意力门，通过对编码阶段传递到解码阶段的特征进行筛选，使模型更集中地利用与目标结构相关的信息 \cite{oktay2018attentionunet}。这类方法之所以重要，不只是因为它提高了目标响应，更因为它减少了解码过程中无关区域带来的干扰。

除局部注意力外，近年来越来越多的研究开始关注更大范围内的特征依赖关系，尤其是自注意力和类 Transformer 结构在分割任务中的应用。其共同特点在于，不再把特征交互严格限制在局部卷积邻域内，而是尝试在更大空间范围内建立关联。对于视网膜血管分割而言，这种思路具有明显吸引力，因为细长血管的走向、分叉与连通关系往往不能仅靠局部纹理完全刻画。也正因为如此，结合卷积与注意力机制的混合结构逐渐增多，其关注点往往集中在高层语义空间中如何更有效地整合全局信息，同时不过度牺牲局部细节 \cite{luo2025panet,zhu2021deformabledetr,chen2024transunet,cao2021swinunet}。

与注意力机制密切相关的另一类方法，是特征融合。编码器—解码器结构虽然天然具备跨层信息传递能力，但如何把浅层细节与深层语义更有效地结合，并不是一个已经完全解决的问题。标准 U-Net 主要依赖跳跃连接完成这一过程，而后续研究则进一步关注不同层级特征之间是否存在语义差距，以及这种差距会不会影响解码阶段的结果恢复。UNet++ 就是这一方向中的代表性工作，它通过嵌套式跳跃连接重构编码器与解码器之间的连接方式，从而缩小不同层级特征之间的语义差异 \cite{zhou2018unetpp}。这一思路说明，特征融合并不仅仅是把多层特征简单拼接起来，更关键的是如何提高融合后的表达质量。

对于视网膜血管分割来说，特征融合的重要性尤其体现在微细血管和复杂分叉区域的恢复上。深层特征通常携带更稳定的语义信息，但浅层特征保留了更多边界、纹理和局部几何细节。如果融合方式过于简单，模型虽然可能获得较强的整体语义判断能力，却未必能在细小结构处维持足够的分辨力。也正因为如此，后续不少方法不再满足于固定形式的跨层连接，而是尝试通过加权、筛选或自适应重组的方式，使不同层级特征以更有针对性的形式参与解码。

在与本文后续工作更直接相关的方法中，PA-Net 提出了 Lightweight Parallel Transformer (LPT) 和 Adaptive Vascular Feature Fusion Module (AFFM) \cite{luo2025panet}。其中，LPT 的重点在于把局部特征提取与较大范围内的依赖建模结合起来，用于增强高层语义空间中的结构表达；AFFM 则着眼于解码阶段的多层特征融合，希望通过更有选择性的方式减少下采样带来的血管信息损失。前者更接近注意力与高层特征建模问题，后者则更直接对应特征融合问题。PA-Net 整体结构及其相关模块示意如图 \ref{fig:panet_lpt_affm_cn} 所示。

\begin{figure}[htbp]
    \centering
    \begin{minipage}{0.4\textwidth}
        \centering
        \includegraphics[width=\linewidth,height=0.22\textheight,keepaspectratio]{pic/panet.png}\\
        (a) PA-Net 整体结构
    \end{minipage}
    \hfill
    \begin{minipage}{0.18\textwidth}
        \centering
        \includegraphics[width=\linewidth,height=0.22\textheight,keepaspectratio]{pic/LPT.png}\\
        (b) LPT 模块
    \end{minipage}
    \hfill
    \begin{minipage}{0.45\textwidth}
        \centering
        \includegraphics[width=\linewidth,height=0.22\textheight,keepaspectratio]{pic/AFFM.png}\\
        (c) AFFM 模块
    \end{minipage}
    \caption{PA-Net 及其与本文后续改进密切相关的两个典型模块示意图。PA-Net 整体结构展示了 LPT 与 AFFM 在网络中的位置；其中，LPT 用于增强高层语义空间中的结构建模能力，AFFM 用于提高解码阶段多层特征融合的有效性。图片来源：文献 \cite{luo2025panet}。}
    \label{fig:panet_lpt_affm_cn}
\end{figure}

从后续研究的角度看，注意力机制与特征融合并不是彼此分离的两条支线，而是共同构成了后文模型改进的重要方法基础。

\subsection{可变形采样与多尺度建模方法}
在标准卷积或固定感受野建模中，特征提取通常围绕规则网格展开，采样位置由卷积核形状预先确定。这种方式在处理一般目标时已经相当有效，但对于细长、弯曲且方向变化明显的结构来说，固定采样方式往往不够灵活。视网膜血管正属于这一类目标：它既不是规则块状区域，也不是边界清晰的刚性结构，而是在空间中不断延伸、分叉并伴随局部宽度变化。对这类目标来说，采样过程如果始终停留在预设位置上，往往很难真正贴近结构本身的走向。

可变形采样方法的基本思想，是让模型不再局限于预先定义好的规则采样点，而是通过学习得到偏移量，使特征聚合能够围绕目标结构自适应展开。以可变形卷积和可变形注意力为代表的方法，都体现了这一思路 \cite{zhu2019dcnv2,zhu2021deformabledetr}。如果把参考位置记为 \(p\)，第 \(k\) 个采样偏移记为 \(\Delta p_k\)，对应权重记为 \(w_k\)，那么这种聚合过程可以概括写成
\[
y(p)=\sum_{k=1}^{K} w_k\,x(p+\Delta p_k),
\]
其中，\(x(\cdot)\) 表示输入特征，\(K\) 表示参与聚合的采样点数。与标准卷积或固定注意力窗口相比，可变形采样更强调“从哪里采样”这一问题，也就是通过动态调整采样位置，让模型能够在弯曲边界、细长分支和复杂分叉区域附近获得更具针对性的特征响应。对视网膜血管分割而言，这种机制之所以重要，正在于血管的连续走向往往无法被规则网格上的少量固定位置充分刻画。

从方法发展来看，可变形卷积更强调在局部邻域中学习更灵活的采样偏移，而可变形注意力则进一步把这种思想推广到更广的特征交互范围。前者通常用于增强卷积操作对不规则几何形态的适应能力，后者则更多用于在特征聚合过程中动态决定不同采样位置的贡献。无论具体形式如何变化，其共同目标都在于减少固定采样方式带来的结构失配，使特征表示能够更贴近目标本身的空间分布。图 \ref{fig:deformable_attention_theory_cn} 给出了固定采样与可变形采样思路的示意对比。

\begin{figure}[htbp]
    \centering
    \includegraphics[width=\textwidth]{pic/Deformable_Attention_fundus.png}
    \caption{眼底血管场景下刚性采样与可变形采样思路的示意对比。左侧表示固定轴向采样可能难以贴合弯曲血管的连通结构，右侧表示自适应采样位置能够更好地沿血管轮廓与分支展开。作者根据可变形采样相关原理绘制 \cite{zhu2019dcnv2,zhu2021deformabledetr}。}
    \label{fig:deformable_attention_theory_cn}
\end{figure}

可变形采样强调的是采样位置的灵活性，而多尺度建模关注的则是不同尺度的信息如何更充分地参与表达。在视网膜血管分割中，粗大血管主干、细小末梢以及复杂分叉区域往往对应不同层级的特征。深层特征通常更有利于形成稳定的全局判断，但浅层特征保留了更多局部边界与纹理细节。单一尺度的特征表达很难同时兼顾这两类信息，因此，多尺度建模逐渐成为增强血管分割能力的重要方向。若将第 \(l\) 个尺度上的输入特征记为 \(x_l\)，对应采样偏移记为 \(\Delta p_{l,k}\)，聚合权重记为 \(w_{l,k}\)，则多尺度可变形聚合可以进一步写为
\[
y(p)=\sum_{l=1}^{L}\sum_{k=1}^{K} w_{l,k}\,x_l(p+\Delta p_{l,k}),
\]
其中，\(L\) 表示参与聚合的尺度数，\(K\) 表示每个尺度上的采样点数。与前面的单尺度形式相比，这一表达强调的是不同层级特征在同一参考位置附近的联合参与方式。

多尺度建模并不只是简单地把不同分辨率的特征同时送入网络，而是强调如何建立不同尺度之间的互补关系。在编码器—解码器结构中，这种互补关系通常体现为跨层特征传递与融合；而在更进一步的方法中，多尺度建模则可能表现为跨层采样、跨尺度聚合，或者在高层表示中主动引入来自浅层特征的补充信息。对细长结构分割任务来说，这一点尤其重要，因为目标既需要全局走向信息来维持连通性，也需要细粒度局部信息来避免末梢血管和复杂分叉处的细节损失。

在与本文后续工作更密切相关的研究中，可变形采样与多尺度建模往往并不是各自独立存在的。相反，它们经常被结合起来使用：一方面通过自适应偏移提高采样位置的灵活性，另一方面通过跨层特征访问与聚合补充不同分辨率下的信息。这种组合思路的优势在于，它既能够增强模型对细长结构实际几何走向的适应能力，也能够缓解高层语义表示中过早丢失浅层细节的问题。后续围绕可变形注意力、调制机制和多尺度扩展展开的改进，本质上都建立在这一方法背景之上。

\subsection{面向细长结构的拓扑约束方法}
对于视网膜血管分割这类细长结构任务来说，像素级分类正确并不一定意味着结构层面也是正确的。血管网络具有明显的连通性、分叉性和延展性，如果模型只是在局部区域给出较高响应，却无法在整体上维持连续走向，最终结果仍可能出现断裂、粘连或孤立噪声区域。也正因为如此，随着分割研究不断深入，越来越多的方法开始不再只盯着局部判别精度，而是进一步关注结果在几何形态和拓扑结构上的合理性。

从结构建模角度看，动态蛇形卷积（Dynamic Snake Convolution, DSConv）代表了一类更强调连续几何约束的思路 \cite{qi2023dsc}。与规则卷积或离散点式采样不同，这类方法试图让相邻采样位置之间形成连续展开的轨迹关系，使采样过程不再只是围绕若干独立位置进行聚合，而是更贴近细长目标在空间中的真实延伸方式。对于血管这类弯曲、分叉且局部宽度不断变化的结构来说，这种建模方式之所以有吸引力，正在于它更容易在特征提取阶段保留目标的连通性和方向性，而不是只强化局部响应。

更具体地说，DSConv 的关键不在于单个采样点是否可变，而在于后续采样位置会依赖前面位置的偏移结果。若将参考位置记为 \(p\)，第 \(k\) 个采样点在水平方向上的基础步长记为 \(\Delta x_k\)，纵向增量记为 \(\delta y_k\)，则蛇形轨迹上的采样位置可以概括写为
\[
\tilde{p}_k = p + \left(\Delta x_k,\; \sum_{t=1}^{k}\delta y_t\right).
\]
这一表达强调了一个重要区别：与离散点式采样中“每个采样点彼此独立”不同，DSConv 中后续位置会沿着前面位置逐步展开，因此更容易形成连续延伸的轨迹。若进一步将沿蛇形轨迹的局部聚合写成
\[
y(p)=\sum_{k=1}^{K} w_k\,x(\tilde{p}_k),
\]
则可以看出，DSConv 的核心是让采样点在几何上更接近细长目标的真实走向，再在这一连续轨迹上完成特征提取与聚合。也正因为如此，它特别适合处理血管、道路、神经纤维等具有较强拓扑连续性的目标。

除了在特征提取阶段引入几何约束，另一类常见做法是在输出端对分割结果进行后处理。对于视网膜血管分割而言，模型输出中常会出现一些面积较小、空间上相互孤立的连通域，它们可能来自背景噪声、病灶反光或局部误检。后处理方法通常不改变网络主体，而是根据结果图中的连通区域大小、形态或连通关系，对明显不合理的预测区域进行筛除或修正。这类方法的思路虽然相对简单，却直接面向结果层面的拓扑整洁性，因此在实际分割流程中具有较强实用性。对细长结构任务来说，后处理的价值并不只是“去噪”，更在于尽量减少那些与目标整体连通形态不一致的局部预测。

与结构约束和输出修正相对应的，还有一类从损失函数层面引入拓扑信息的方法。clDice 就是其中具有代表性的工作，它试图通过骨架层面的重叠关系，把连通性约束直接纳入优化目标 \cite{shit2021cldice}。若将预测结果记为 \(P\)，真实标注记为 \(G\)，则 clDice 可以写为
\[
\mathrm{clDice}(P,G)=
\frac{2\,T_{\mathrm{prec}}(P,G)\,T_{\mathrm{sens}}(P,G)}
{T_{\mathrm{prec}}(P,G)+T_{\mathrm{sens}}(P,G)},
\]
其中，\(T_{\mathrm{prec}}\) 和 \(T_{\mathrm{sens}}\) 分别表示基于骨架重叠关系定义的拓扑精确率和拓扑灵敏度。与只关注区域重叠程度的传统损失不同，这类方法更加关心预测结果和真实结构在中心线或骨架层面的对应关系，因此更适合处理血管、神经纤维、道路等细长目标。对视网膜血管分割而言，这一思路的意义在于，它不只是鼓励模型在像素层面接近标注结果，还希望模型在更高层次上保留结构的连续性。

从已有方法来看，面向细长结构的拓扑约束大致可以从几个角度来理解：有些工作在特征提取阶段引入更贴近几何轨迹的建模方式，有些工作在网络输出后利用连通域规则对结果进行修正，也有一些工作在训练阶段通过附加损失对结构连续性施加约束。它们关注的层面并不相同，前者更接近结构表达问题，中间一种更偏向结果清理，后一种则更强调优化目标的引导作用。对本文而言，这些方法主要作为后续改进思路的背景参照，并不意味着在当前实验条件下都需要被同等展开。它们共同指向的，仍然是如何让分割结果不仅“分得出来”，而且在结构上尽量接近真实血管网络。

对本文后续工作而言，这类方法背景的重要性在于，它说明了结构连续性并不是一个附带问题，而是值得单独处理的关键因素。后文围绕连续几何建模、后处理以及拓扑相关损失展开的改进，正是在这一方法脉络下逐步展开的。

\subsection{本章小结}
本章围绕本文后续模型改进所涉及的理论基础与相关方法进行了梳理。从视网膜血管分割任务本身出发，本章说明了该任务在输入输出形式、数据集差异、评价指标以及结构连续性要求等方面的基本特点；接着介绍了 U-Net 与 nnU-Net 的基本原理，明确了编码—解码结构、跳跃连接以及自适应配置框架在医学图像分割中的基础作用；随后又进一步梳理了注意力机制、特征融合、可变形采样、多尺度建模以及面向细长结构的拓扑约束方法。整体来看，这些内容分别从信息选择、跨层表达、采样位置组织和结构连续性保持等角度，为后续模型设计提供了方法基础。

这一章并不直接讨论具体实验结果，而是为后文的设计与实现提供背景依据。也正因为前面已经把相关方法脉络梳理清楚，第三章中围绕基线建立、结构迁移、主线改进以及结果优化展开的设计过程，才会有更明确的逻辑起点。

\section{改进模型设计与实现}

\subsection{整体研究思路与实验框架}
为了使后续各阶段模型比较建立在统一基础之上，本文在展开具体结构改进之前，先搭建了较为清晰的实验框架。整体研究思路不是在不同设置下零散尝试各种模块，而是在前期探索的基础上逐步收敛到统一协议，再在同一实验条件下推进后续结构改进与结果优化。按研究推进顺序来看，全文实验大体可以分为两个阶段：前一阶段主要用于建立初始参照并判断常规训练层面调整的作用，后一阶段则在统一数据条件和训练协议下，对后续各类结构改动进行系统比较。

前期探索阶段对应最早的几组对照实验，重点在于建立初始基线，并观察输入增强和损失函数调整对分割性能的影响。该阶段的训练与验证基于多数据源样本，主要涉及 STARE、CHASE\_DB1 和 HRF 等公开眼底数据 \cite{hoover2000locating,chasedb12013,budai2013robust}，并沿用既有划分方式完成训练与验证；与此同时，DRIVE 数据集作为测试集，用于在相对统一的目标域上观察模型的初始表现。为了尽量保留彩色眼底图像中的原始信息，前期输入采用三通道分离形式，即先将原始图像转换为 nnU-Net 可读取的 NIfTI 格式，再将红、绿、蓝三个通道分别作为输入。这样的安排，更多是为后续正式比较之前提供一个相对宽一点的前期参照，而不是直接形成最终方案。

后期主实验则从统一协议建立开始，研究重心也由前期可行性判断逐步转向结构改进。该阶段将数据来源固定在同一公开基准（DRIVE）上，并采用一致的数据组织方式和五折交叉验证协议，以减少单一划分带来的影响。与前期阶段不同，后期主实验不再使用三通道彩色输入，而是采用经过增强预处理的单通道图像。具体而言，原始眼底图像依次经过灰度化、标准化、自适应直方图均衡化和 gamma 校正，再转换为单通道 NIfTI 文件输入模型。该预处理流程参考了 PA-Net 的相关做法 \cite{luo2025panet}，目的是增强细小血管和低对比度区域的可分辨性，并尽量减少无关颜色信息的干扰。自后期基线建立之后，后续各组结构改进实验均在这一统一实验基准下展开，这样可以尽可能控制数据条件变化带来的影响，使不同结构之间的横向比较建立在同一基础之上。

在训练设置上，全部实验均基于 nnU-Net v2 框架实现 \cite{isensee2021nnunet}。前期探索阶段主要用于快速建立参照和观察趋势，因此沿用固定划分，并采用相对简化的对比方式；后期主实验则统一采用 nnU-Net v2 的 2D 配置，输入 patch size 设为 \(512 \times 512\)。多数结构实验训练 100 轮，基线模型的 trainer 主要遵循 nnU-Net 默认优化设置，而 PA-Net 相关结构变体则使用迁移模块所需的自定义 trainer。因此，后文分析主要把共同的数据划分、输入组织方式、推理流程和评价脚本作为比较基础，避免把性能变化简单归因到某一个孤立超参数上。

表 \ref{tab:protocol_overview_cn} 对两个实验阶段的差异作了简要概括。

\begin{table}[htbp]
\centering
\caption{两个实验阶段的整体对比}
\label{tab:protocol_overview_cn}
\resizebox{\textwidth}{!}{
\begin{tabular}{llll}
\hline
阶段 & 输入与比较设置 & 训练设置 & 主要作用 \\
\hline
前期探索 & 多来源公开眼底样本，三通道 NIfTI 输入 & 固定划分；前期设置中主要训练 150 轮 & 建立初始参照，观察输入增强和损失调整作用 \\
后期主实验 & 统一公开基准，增强单通道 NIfTI 输入 & 五折验证；多数结构实验训练 100 轮 & 比较结构变体并完成最终结果优化 \\
\hline
\end{tabular}
}
\end{table}

在评价方式上，前期探索阶段主要采用 IoU、mIoU、Recall、Precision、Accuracy、mPA 和 F1-score 等指标，用于建立初始比较参照；后期主实验则统一采用 Sensitivity (SE)、Specificity (SP)、Accuracy (ACC)、Precision (PR) 和 F1-score 进行评估。考虑到当前任务更关注细小血管恢复与误检控制，后文分析会更重视 SE、PR 和 F1-score 的变化，同时结合 SP 与 ACC 描述整体预测平衡。对于五折交叉验证结果，本文不直接依据单个折作出结论，而是对所有折的预测结果进行统一汇总后再计算指标，以减少单次划分带来的波动，并提高不同模型之间比较的一致性。表中带有 ``\(\pm\)'' 的结果表示折层面的均值和标准差；只报告单一数值的结果，则对应最终用于比较的汇总预测结果。

\subsection{基线模型与前期探索}
在完成整体实验框架的建立之后，本文首先以前期基线实验以及围绕基线展开的输入增强、损失函数探索作为后续研究的起点。这个阶段并不以直接形成最终方案为目标，而是希望先回答一个更基础的问题：在 nnU-Net 已有内部预处理与训练机制的基础上，额外的输入端增强和损失函数调整究竟还能带来多大程度的改善。由于这一阶段仍属于前期可行性判断，实验沿用既有训练集和验证集划分，并统一以 DRIVE 作为测试集，在相同数据条件下比较不同设置的表现。

第一次实验采用标准 nnU-Net 基线模型 \cite{isensee2021nnunet}，作为本文最初的性能参照。该实验使用原始输入形式和常规训练设置进行训练，主要是为了观察在不引入额外结构改动的情况下，当前框架本身能够达到怎样的分割效果。结果提示，标准 nnU-Net 在这一阶段已经能够取得相对稳定的整体表现，因此可以作为后续探索的起点。

第二次实验在相同训练框架下引入 CLAHE 图像增强，用于观察输入增强对血管分割的影响。这样做的考虑在于，眼底图像中的细小血管和低对比度区域往往不够明显，适度增强局部对比度可能有助于改善模型对这些区域的感知。图 \ref{fig:clahe_compare_cn} 给出了原始眼底图像与 CLAHE 预处理结果的直观对比。可以看到，经过 CLAHE 处理后，局部对比度有所增强，部分细小血管区域变得更加清晰，但同时也会增强部分非血管纹理与背景响应。实验结果显示，CLAHE 对训练初期的收敛过程有一定帮助，但对最终分割性能的提升有限。对这一组前期实验而言，更合适的解释并不是这类处理完全没有作用，而是 nnU-Net 自身已经具备较为完善的内部预处理流程，因此在图像进入网络前再做一轮额外增强，并没有带来预期中的明显增益。

\begin{figure}[htbp]
    \centering
    \begin{minipage}{0.42\textwidth}
        \centering
        \includegraphics[width=\textwidth]{pic/original_image.jpg}\\
        (a) 原始眼底图像
    \end{minipage}
    \hfill
    \begin{minipage}{0.42\textwidth}
        \centering
        \includegraphics[width=\textwidth]{pic/CLAHE_image.jpg}\\
        (b) CLAHE 预处理结果
    \end{minipage}
    \caption{原始眼底图像与 CLAHE 预处理结果对比。经过 CLAHE 处理后，图像局部对比度有所增强，部分细小血管区域变得更加清晰，但同时也会增强部分非血管纹理与背景响应。}
    \label{fig:clahe_compare_cn}
\end{figure}

第三次实验在 CLAHE 的基础上进一步引入 Focal Loss，以考察损失函数调整对困难样本和类别不平衡问题的作用。与基线和单纯输入增强相比，这一设置使模型收敛更快，但最终性能并未继续提升。在当前结构保持不变的前提下，损失函数层面的调整仍更多作用于训练过程本身，而没有带来稳定的结果改观。

为便于对前期探索结果形成统一参照，表 \ref{tab:early_experiments_cn} 对前三次实验的设置与代表性结果进行了汇总。

\begin{table}[htbp]
\centering
\caption{前期探索实验结果汇总}
\label{tab:early_experiments_cn}
\resizebox{\textwidth}{!}{%
\begin{tabular}{llllll}
\hline
实验编号 & 方法 & 训练设置 & 收敛情况 & 代表性结果 & 主要观察 \\
\hline
Exp1 & nnU-Net 基线 & Dice, 150 epoch & 第69个 epoch 收敛 & F1 = 0.7924 & 建立前期基线 \\
Exp2 & CLAHE + nnU-Net & Dice, 150 epoch & 第60个 epoch 收敛 & F1 = 0.7827 & 改善早期收敛，但最终提升有限 \\
Exp3 & CLAHE + nnU-Net + Focal Loss & Focal Loss, 150 epoch & 第53个 epoch 收敛 & F1 = 0.7802 & 收敛更快，但性能未进一步提升 \\
\hline
\end{tabular}
}
\end{table}

回看这一阶段可以发现，输入增强和损失函数调整会影响训练过程，也会带来局部指标上的变化，但这种影响更多停留在输入表征与优化目标层面，并未直接解决细小血管恢复不足和结构连续性建模不充分的问题。也正因为如此，后续研究的重点才逐步从训练层面的调整转向网络结构本身的改进。

\subsection{PA-Net 模块迁移与消融实验}
在前期探索表明额外输入端预处理并未带来明显收益、损失函数调整也未形成稳定增益之后，本文开始把研究重点转向网络结构本身。作为这一阶段的起点，本文首先参考 PA-Net 的设计思路，尝试将其中两个与视网膜血管分割关系较为密切的核心模块迁移到 nnU-Net 框架中，即 Lightweight Parallel Transformer (LPT) 和 Adaptive Vascular Feature Fusion Module (AFFM) \cite{luo2025panet}。之所以从这两个模块入手，是因为它们分别对应了后续研究最关心的两个问题：一是高层特征空间中的结构表达能力，二是解码阶段不同层级特征的有效融合。

从方法定位上看，LPT 被置于瓶颈层（bottleneck）位置，主要用于增强模型对长程依赖和全局结构信息的建模能力；AFFM 则用于替代标准 skip connection 的融合方式，希望通过自适应加权减轻下采样带来的细节损失。由于第二章已经对相关原理作了介绍，这里不再展开模块结构本身，而是重点考察这些模块迁移到 nnU-Net 之后能否在统一协议下产生稳定收益。换句话说，这一阶段的任务并不是立即得到最终方案，而是先判断 PA-Net 的结构思路在当前框架中到底能不能成立。

基于这一考虑，本文首先构建完整 PA-Net 迁移模型，并在同一实验设置下进一步拆分出两个单模块版本，即仅保留 LPT 的 LPT Only 和仅保留 AFFM 的 AFFM Only，用于判断完整迁移结果中的性能变化主要来自哪个模块。相关结果见表 \ref{tab:lpt_affm_ablation_cn}。

\begin{table}[htbp]
\centering
\caption{PA-Net 模块迁移与消融实验结果对比}
\label{tab:lpt_affm_ablation_cn}
\resizebox{\textwidth}{!}{
\begin{tabular}{lccccc}
\hline
实验组 & SE & SP & ACC & PR & F1-score \\
\hline
nnU-Net v2 Baseline & $0.7627 \pm 0.0446$ & $0.9852 \pm 0.0050$ & $0.9659 \pm 0.0045$ & $0.8238 \pm 0.0646$ & $0.7891 \pm 0.0320$ \\
PA-Net (Full) & $0.7211 \pm 0.0763$ & $0.9888 \pm 0.0050$ & $0.9659 \pm 0.0059$ & $0.8570 \pm 0.0621$ & $0.7782 \pm 0.0464$ \\
LPT Only & $0.7597 \pm 0.0492$ & $0.9864 \pm 0.0041$ & $0.9669 \pm 0.0044$ & $0.8361 \pm 0.0546$ & $0.7935 \pm 0.0297$ \\
AFFM Only & $0.7573 \pm 0.0475$ & $0.9868 \pm 0.0039$ & $0.9672 \pm 0.0045$ & $0.8400 \pm 0.0517$ & $0.7942 \pm 0.0300$ \\
\hline
\end{tabular}
}
\end{table}

从表 \ref{tab:lpt_affm_ablation_cn} 可以看出，完整 PA-Net 迁移结果并未优于基线。虽然 PA-Net (Full) 在 PR 上有所提高，提示预测更为保守，但其 SE 明显下降，最终 F1-score 也低于基线。这意味着完整迁移后的模型更倾向于保留高置信度血管区域，却牺牲了部分细小血管和弱响应区域的检出能力。相比之下，两组单模块消融的结果整体上更为稳定。LPT Only 在基本保持 SE 的同时略微提高了 PR 和 F1-score，提示在瓶颈层引入更强的高层结构建模具有一定参考价值；AFFM Only 虽然取得了本组实验中较高的 ACC 和 F1-score，但它相对 LPT Only 的 F1-score 优势非常有限，不能单独说明后续主线应转向解码端融合。因此，本文在下一步继续通过深度监督控制实验检查 AFFM 的稳定性，而 LPT 仍作为后续瓶颈层采样改进路线中更相关的分支。

在这一结果基础上，本文进一步把问题缩小到 AFFM 本身。作为解码阶段的多层特征融合模块，AFFM 的收益是否会受到 nnU-Net 默认深度监督机制的影响，是这一阶段需要进一步确认的。考虑到深度监督同样作用于解码器的多个输出层，并通过辅助分割头强化多尺度特征学习 \cite{isensee2021nnunet,lee2015dsn}，本文在 AFFM Only 的基础上关闭深度监督，构造 AFFM w/o DS 模型，以考察二者之间是否存在潜在冲突。若将主输出损失记为 \(\mathcal{L}_{\mathrm{main}}\)，将第 \(s\) 个辅助输出损失记为 \(\mathcal{L}_{\mathrm{aux}}^{(s)}\)，则深度监督下的总损失可以写为
\[
\mathcal{L} = \mathcal{L}_{\mathrm{main}} + \sum_{s=1}^{S}\lambda_s \mathcal{L}_{\mathrm{aux}}^{(s)},
\]
其中，\(\lambda_s\) 表示不同辅助分支的权重。其训练机制示意如图 \ref{fig:deep_supervision_cn_33} 所示。

\begin{figure}[htbp]
\centering
\includegraphics[width=0.95\textwidth]{pic/deep_supervision.png}
\caption{nnU-Net 解码器中深度监督机制示意图。训练阶段在多个解码层附加辅助分割头并共同计算加权损失，推理阶段仅保留主输出。}
\label{fig:deep_supervision_cn_33}
\end{figure}

相应实验结果如表 \ref{tab:affm_ds_cn_33} 所示。

\begin{table}[htbp]
\centering
\caption{AFFM 在有无深度监督下的结果对比}
\label{tab:affm_ds_cn_33}
\resizebox{\textwidth}{!}{
\begin{tabular}{lccccc}
\hline
实验组 & SE & SP & ACC & PR & F1-score \\
\hline
AFFM Only (with DS) & $0.7573 \pm 0.0475$ & $0.9868 \pm 0.0039$ & $0.9672 \pm 0.0045$ & $0.8400 \pm 0.0517$ & $0.7942 \pm 0.0300$ \\
AFFM w/o DS & $0.7401 \pm 0.0493$ & $0.9880 \pm 0.0038$ & $0.9667 \pm 0.0048$ & $0.8490 \pm 0.0534$ & $0.7883 \pm 0.0308$ \\
\hline
\end{tabular}
}
\end{table}

从表 \ref{tab:affm_ds_cn_33} 可以看出，关闭深度监督后，AFFM 模型的 SP 和 PR 略有提高，但 SE 和 F1-score 均有所下降，整体上呈现出更保守的预测倾向。去除深度监督后，模型虽然减少了部分伪阳性，但对细小血管和弱响应区域的检出能力也随之减弱。因此，仅靠关闭深度监督并不能改善 AFFM 的整体表现，前面观察到的性能变化也不应只归因于“深度监督冲突”。

综合这一节的结果可以看出，PA-Net 中单个模块的设计思想具有一定参考价值，但直接将完整结构迁移到 nnU-Net 中并不能形成更优方案。进一步地，AFFM 虽然在局部比较中表现出一定潜力，但其收益并不稳定，且在与深度监督关系的控制实验中也未显示出更明确的优势。相比之下，LPT 在这一阶段表现出更清晰的继续探索价值。因此，这一节的意义不在于确定最终模型，而在于完成第一轮结构筛选：保留值得继续推进的方向，并为下一步围绕 LPT 展开的结构改进奠定基础。

\subsection{基于 LPT 的结构改进方法}
在上一节完成第一轮结构筛选之后，本文的后续改进主线逐步收拢到 LPT 模块。原因也比较明确：与完整 PA-Net 迁移和 AFFM 相关探索相比，LPT 在当前框架下表现出更稳定的继续优化价值。前一阶段的结论并不是“PA-Net 无效”，而是提示本文后续更值得保留的是其中与高层结构建模相关的部分。因此，本节不再继续围绕解码端融合机制展开，而是把关注点重新放回瓶颈层位置的 LPT，并围绕其采样方式和特征聚合方式展开后续改进。

这一轮改进首先从可变形采样开始。本文参考可变形注意力 \cite{zhu2021deformabledetr} 的思路，将 PA-Net 中 LPT 原有的 Gated axial-attention（GAA，门控轴向注意力）操作替换为更灵活的可变形采样机制。在 PA-Net 中，GAA 主要沿宽度和高度两个方向进行门控多头注意力建模，并与局部卷积分支结合 \cite{luo2025panet}。这种设计有助于在两个主要空间方向上组织特征交互，但当血管强烈弯曲、出现复杂分叉或非常细小时，其空间交互方式仍然不够灵活。相较之下，可变形采样允许采样位置围绕参考点进行自适应偏移，使模型在高层语义空间中有更多机会沿着弯曲血管轨迹聚合特征。由于第二章已经对固定采样与可变形采样的原理差异进行了说明，这里不再重复展开方法定义，而是直接关注这种替换在实验中的实际影响。

围绕这一步替换，本文先后构建了三个对照模型：其一为去除 AFFM 后的 Deformable LPT（4 point），用于判断可变形采样替换本身是否有效；其二是在相同结构下重新加入 AFFM，用于确认在新的 LPT 形式下是否仍有保留 AFFM 的必要；其三则将采样点数由 4 point 扩展到 8 point，用于观察采样规模变化对结果的影响。相关结果如表 \ref{tab:deformable_lpt_cn_34} 所示。

\begin{table}[htbp]
\centering
\caption{可变形 LPT 相关实验结果对比}
\label{tab:deformable_lpt_cn_34}
\resizebox{\textwidth}{!}{
\begin{tabular}{lccccc}
\hline
实验组 & SE & SP & ACC & PR & F1-score \\
\hline
LPT Only & $0.7597 \pm 0.0492$ & $0.9864 \pm 0.0041$ & $0.9669 \pm 0.0044$ & $0.8361 \pm 0.0546$ & $0.7935 \pm 0.0297$ \\
Deformable LPT (No AFFM, 4 point) & $0.7634 \pm 0.0540$ & $0.9863 \pm 0.0043$ & $0.9672 \pm 0.0044$ & $0.8371 \pm 0.0511$ & $0.7958 \pm 0.0275$ \\
Deformable LPT + AFFM & $0.7571 \pm 0.0495$ & $0.9867 \pm 0.0042$ & $0.9670 \pm 0.0044$ & $0.8388 \pm 0.0545$ & $0.7933 \pm 0.0301$ \\
Deformable LPT (No AFFM, 8 point) & $0.7573 \pm 0.0506$ & $0.9869 \pm 0.0038$ & $0.9672 \pm 0.0047$ & $0.8412 \pm 0.0504$ & $0.7946 \pm 0.0296$ \\
\hline
\end{tabular}
}
\end{table}

从表 \ref{tab:deformable_lpt_cn_34} 可以看出，在不引入 AFFM 的条件下，将 LPT 替换为可变形采样后，模型的 SE 和 F1-score 较 LPT Only 有所提高，其中 4 point 设置取得了本组实验中的最高 F1-score。SP 和 ACC 与 LPT Only 接近，提示这一变化并不是以明显牺牲背景抑制为代价。整体来看，表中结果支持在下一轮实验中继续保留更灵活的瓶颈层采样机制。与此同时，重新加入 AFFM 并未带来进一步提升，F1-score 反而从 0.7958 降至 0.7933，这进一步提示 AFFM 在当前框架下未表现出稳定增益。至于 8 point 设置，虽然在 PR 上略有提高，但 SE 和 F1-score 均未超过 4 point 版本，因此后续实验将不含 AFFM 的 Deformable LPT（4 point）保留为新的改进基线。

在此基础上，本文继续围绕这一改进基线向前推进。首先，在单尺度可变形采样框架内进一步加入调制权重（modulation mask），用于调节不同采样点的贡献，形成 Modulated Deformable（4 point）；随后，又在这一结构基础上把采样范围从当前层扩展到编码器多层特征，通过跨层访问和聚合补充更多高分辨率信息，形成 MS-Modulated（4 point）。后者可以看作是在 Modulated Deformable 的基础上进一步引入多尺度采样后的扩展形式。图 \ref{fig:modulated_deformable_cn_34} 和图 \ref{fig:ms_modulated_cn_34} 分别给出了这两步改进的示意。

\begin{figure}[htbp]
\centering
\includegraphics[width=0.9\textwidth]{pic/Modulated_Deformable_Sampling_fundus.png}
\caption{单尺度 Modulated Deformable 采样示意图。查询点通过可变形偏移访问周围采样位置，并利用调制权重对不同采样点的贡献进行进一步调节。作者根据本文实现和可变形采样原理绘制 \cite{zhu2019dcnv2,zhu2021deformabledetr}。}
\label{fig:modulated_deformable_cn_34}
\end{figure}

\begin{figure}[htbp]
\centering
\includegraphics[width=0.95\textwidth]{pic/MS_modulated_DF.png}
\caption{多尺度 Modulated Deformable 采样示意图。瓶颈层特征作为 query，跨层访问编码器多尺度特征，并通过加权聚合形成增强后的瓶颈层表示。作者根据本文实现绘制。}
\label{fig:ms_modulated_cn_34}
\end{figure}

在上述递进关系基础上，本文进一步构建了 Modulated Deformable（4 point）和 MS-Modulated（4 point）两个模型，并在统一实验设置下进行比较。相关结果如表 \ref{tab:modulated_msdeform_cn_34} 所示。

\begin{table}[htbp]
\centering
\caption{Modulated 与多尺度可变形 LPT 实验结果对比}
\label{tab:modulated_msdeform_cn_34}
\resizebox{\textwidth}{!}{
\begin{tabular}{lccccc}
\hline
实验组 & SE & SP & ACC & PR & F1-score \\
\hline
Deformable LPT (No AFFM, 4 point) & $0.7634 \pm 0.0540$ & $0.9863 \pm 0.0043$ & $0.9672 \pm 0.0044$ & $0.8371 \pm 0.0511$ & $0.7958 \pm 0.0275$ \\
Deformable LPT (No AFFM, 8 point) & $0.7573 \pm 0.0506$ & $0.9869 \pm 0.0038$ & $0.9672 \pm 0.0047$ & $0.8412 \pm 0.0504$ & $0.7946 \pm 0.0296$ \\
Modulated Deformable (4 point) & $0.7660 \pm 0.0511$ & $0.9858 \pm 0.0048$ & $0.9669 \pm 0.0046$ & $0.8324 \pm 0.0560$ & $0.7950 \pm 0.0305$ \\
MS-Modulated (4 point) & $0.7582 \pm 0.0509$ & $0.9870 \pm 0.0038$ & $0.9673 \pm 0.0046$ & $0.8419 \pm 0.0497$ & $0.7954 \pm 0.0283$ \\
\hline
\end{tabular}
}
\end{table}

从表 \ref{tab:modulated_msdeform_cn_34} 可以看出，Modulated Deformable（4 point）在 SE 上取得了本组实验中的最高值，提示调制权重与更强的细小血管和弱响应区域检出能力相关。在单尺度可变形采样中，额外的选择机制可能改变了采样分布，使模型更倾向于关注与血管连续性相关的位置。在此基础上进一步引入多尺度采样后，MS-Modulated（4 point）并未继续提高 SE，但在 SP、ACC 和 PR 上取得了更好的表现，提示跨层多尺度采样有助于提高模型对背景噪声的抑制能力，并改善整体预测平衡。因此，从 Modulated Deformable 到 MS-Modulated，模型的改进重点由增强检出能力进一步转向不同层级信息的补充与融合。

综合这一节可以看到，围绕 LPT 展开的这一轮结构改进已经形成了较为清晰的递进关系：先通过可变形采样解决固定采样方式对细长血管适配不足的问题，再通过调制机制增强单尺度采样中的关键位置选择能力，最后进一步引入多尺度采样与聚合，以补充不同层级的特征信息。虽然这一阶段各模型之间的性能差异仍然有限，但研究路线已经从“模块能否迁移”进一步转向“如何在统一框架中形成更稳定的结构改进逻辑”。在这一基础上，下一节将继续沿着多尺度建模方向展开，并进一步引入注意力增强与连续几何建模。

\subsection{多尺度建模与注意力增强方法}
在上一节中，本文已经把改进主线推进到 MS-Modulated（4 point）阶段。到这一步为止，模型已经具备了跨层多尺度特征访问与聚合能力，但新的问题也随之显现出来：虽然多尺度采样有助于补充浅层细节，并改善整体预测平衡，但仅有跨层聚合本身可能仍不足以充分强化对细长血管方向和空间位置的选择性表达。基于这一观察，本文首先在 MS-Modulated 结构中进一步引入 Coordinate Attention (CA) \cite{hou2021coordatt}，希望在已有多尺度建模基础上继续增强模型对方向信息和位置信息的建模能力。

与一般通道注意力不同，Coordinate Attention 的特点在于同时保留通道关系与空间方向信息。其基本做法是分别沿水平和垂直方向对特征进行聚合，再通过共享映射和方向相关变换得到两个方向敏感的注意力图。若将输入特征记为 \(X \in \mathbb{R}^{C\times H\times W}\)，则其输出可以概括为
\[
Y_{\mathrm{CA}} = X \odot a_h \odot a_w,
\]
其中，\(a_h\) 和 \(a_w\) 分别表示沿两个空间方向生成的注意力响应。对视网膜血管分割而言，这种设计的意义在于，它并不是简单增强所有通道，而是在已有多尺度聚合结果基础上进一步强调与血管方向、位置和连续走向相关的响应区域。从实现位置上看，CA 被加入到多尺度 LPT 模块的入口处，对瓶颈层特征先进行方向相关的重标定，再送入后续局部特征提取与多尺度可变形聚合过程。图 \ref{fig:ms_ca_cn_35} 给出了 MS-Modulated 与 Coordinate Attention 结合后的结构示意。

\begin{figure}[htbp]
\centering
\includegraphics[width=0.95\textwidth]{pic/MS_CA.png}
\caption{MS-Modulated 与 Coordinate Attention 结合的结构示意图。多尺度可变形采样先完成跨层信息聚合，随后由 Coordinate Attention 在空间方向上进一步重标定特征响应。作者根据本文实现和 Coordinate Attention 原理绘制 \cite{hou2021coordatt}。}
\label{fig:ms_ca_cn_35}
\end{figure}

基于上述思路，本文将 MS-Modulated（4 point）作为直接对照组，并进一步构建 MS-Modulated + CA 模型，以考察在已有跨层多尺度采样基础上引入显式位置注意力后，模型是否能够取得更稳定的综合表现。相关结果如表 \ref{tab:msmod_ca_cn_35} 所示。

\begin{table}[htbp]
\centering
\caption{MS-Modulated 与 MS-Modulated + CA 实验结果对比}
\label{tab:msmod_ca_cn_35}
\resizebox{\textwidth}{!}{
\begin{tabular}{lccccc}
\hline
实验组 & SE & SP & ACC & PR & F1-score \\
\hline
MS-Modulated (4 point) & $0.7582$ & $0.9870$ & $0.9673$ & $0.8419$ & $0.7954$ \\
MS-Modulated + CA & $0.7697$ & $0.9857$ & $0.9673$ & $0.8324$ & $0.7973$ \\
\hline
\end{tabular}
}
\end{table}

从表 \ref{tab:msmod_ca_cn_35} 可以看出，在 MS-Modulated 的基础上引入 CA 后，模型在 SE 和 F1-score 上进一步提高，其中 SE 达到 0.7697，F1-score 提升到 0.7973。这一结果提示，在已有跨层可变形采样基础上，显式位置注意力可能进一步支持细小血管和弱响应区域的检出，并改善整体分割质量。与前一节相比，MS-Modulated 更偏向于提升 SP、ACC 和 PR，而加入 CA 后，模型又在 SE 和 F1-score 上取得了更好的结果。CA 的作用因此不只是继续压低误检，也体现在已有多尺度特征交互基础上强化与血管方向和位置信息相关的表达。

在此基础上，本文进一步关注到另一个问题：虽然多尺度可变形采样与 Coordinate Attention 已经增强了跨层信息交互和空间位置选择能力，但当前采样方式本质上仍更接近离散点式建模，能够找到关键位置，却未必能够充分描述血管这种细长、连续、弯曲的几何形态。基于这一观察，后续实验进一步引入动态蛇形卷积（Dynamic Snake Convolution, DSConv）\cite{qi2023dsc}，希望让采样点之间形成连续轨迹约束，从而更贴合视网膜血管的真实走向。需要明确的是，这一步并不是将整个网络中的普通卷积全部替换为蛇形卷积，而是将其用于瓶颈层 LPT 模块中的多尺度采样分支，用来替代前一阶段 MS-Modulated + CA 中的离散点式跨层可变形采样。

与前面各采样点相对独立的离散采样方式不同，动态蛇形卷积的关键在于让后续采样位置依赖前面位置的偏移结果。若将第 \(k\) 个采样点的纵向增量记为 \(\delta y_k\)，横向基础步长记为 \(\Delta x_k\)，则蛇形轨迹上的采样位置可以概括为
\[
\tilde{p}_k = p + \left(\Delta x_k,\; \sum_{t=1}^{k}\delta y_t\right).
\]
这个式子说明，DSConv 不再把各采样点看作彼此独立的离散位置，而是让它们沿着一条逐步展开的轨迹参与特征提取。对视网膜血管分割而言，这一点很关键，因为血管的弯曲走向、分叉关系和局部宽度变化，本来就更接近连续形态而不是独立点集。

在本文的实现中，这种连续轨迹并不是停留在单尺度局部采样上，而是进一步嵌入到多尺度聚合过程中。若将第 \(l\) 个尺度上的输入特征记为 \(x_l\)，经线性映射后的表示记为 \(P_l(x_l)\)，对应权重记为 \(A_{l,k}\) 和 \(m_{l,k}\)，则多尺度动态蛇形聚合可以概括写为
\[
\mathrm{MSDSConv}(p) = \sum_{l=1}^{L}\sum_{k=1}^{K} A_{l,k} m_{l,k}\,P_l\big(x_l(\tilde{p}_{l,k})\big).
\]
与前一阶段的多尺度调制可变形采样形式相比，这里的主要变化并不只是“继续做多尺度聚合”，而是把原本离散的跨层采样点替换为沿连续轨迹展开的位置序列。这样一来，模型在访问不同层级特征时，不再只是寻找若干独立位置的响应，而是尝试沿着更符合血管几何走向的路径完成聚合。图 \ref{fig:msds_ca_cn_35} 展示了 MS-DSConv 与 Coordinate Attention 结合后的整体结构示意。

\begin{figure}[htbp]
\centering
\includegraphics[width=0.95\textwidth]{pic/MSDS.png}
\caption{MS-DSConv 与 Coordinate Attention 结合的结构示意图。多尺度采样分支中的采样点沿连续蛇形轨迹展开，在完成跨层加权聚合后，再由 Coordinate Attention 对空间响应进行进一步重标定。作者根据本文实现、Dynamic Snake Convolution 和 Coordinate Attention 原理绘制 \cite{qi2023dsc,hou2021coordatt}。}
\label{fig:msds_ca_cn_35}
\end{figure}

基于上述思路，本文将 MS-Modulated + CA 作为直接对照组，并进一步构建 MS-DSConv + CA 模型，以考察在已有多尺度采样和空间注意力基础上引入连续几何轨迹建模后，模型是否能够进一步提升细小血管的恢复能力和整体结构完整性。相关结果如表 \ref{tab:msdsconv_ca_cn_35} 所示。

\begin{table}[htbp]
\centering
\caption{MS-Modulated + CA 与 MS-DSConv + CA 实验结果对比}
\label{tab:msdsconv_ca_cn_35}
\resizebox{\textwidth}{!}{
\begin{tabular}{lccccc}
\hline
实验组 & SE & SP & ACC & PR & F1-score \\
\hline
MS-Modulated + CA & $0.7697$ & $0.9857$ & $0.9673$ & $0.8324$ & $0.7973$ \\
MS-DSConv + CA & $0.8083$ & $0.9815$ & $0.9666$ & $0.7997$ & $0.8020$ \\
\hline
\end{tabular}
}
\end{table}

从表 \ref{tab:msdsconv_ca_cn_35} 可以看出，在 MS-Modulated + CA 的基础上进一步引入动态蛇形卷积后，模型在 SE 和 F1-score 上取得提高，其中 SE 提升到 0.8083，F1-score 提升到 0.8020。这一结果与 DSConv 的设计目标基本一致：当采样机制从离散点式聚合过渡到连续轨迹式建模后，模型更倾向于恢复细小血管、断裂区域以及连续拓扑结构。与此同时，MS-DSConv + CA 的 SP 和 PR 相比前一阶段略有下降，其中 PR 为 0.7997，低于 0.8。更强调几何连续性的建模方式可能会使模型更愿意标记潜在血管像素，从而增加一部分伪阳性和小噪声区域。

综合这一节的结果可以看出，在多尺度建模主线上继续引入显式位置注意力与连续几何轨迹约束是值得保留的方向。其中，MS-Modulated + CA 在保持整体平衡的同时提高了细小血管检出能力，而 MS-DSConv + CA 则在结构连续性建模上继续向前推进，并取得了这一阶段最高的 F1-score。与此同时，其偏低的 PR 也提示结果中仍存在一定误检区域。也正因为如此，下一节将继续从输出结果优化和损失函数约束两个层面展开补充探索，以检查最终输出是否能够变得更干净、更稳定。

\subsection{后处理与损失函数探索}
在上一节中，MS-DSConv + CA 取得了当前阶段最高的 F1-score，并带来了更高的 SE，这提示多尺度建模、位置注意力和连续几何轨迹约束组成的主线对于细小血管和结构连续性建模是有帮助的。然而，新的问题也随之显现出来：该模型虽然在 SE 和 F1-score 上表现较好，但其 PR 仍低于 0.8。这意味着模型在更积极地恢复细小血管和连续结构的同时，也保留了一部分小噪声连通域和伪阳性区域。于是，这一阶段的后续探索开始转向两个更直接的问题：一是如何从输出端清理明显不合理的小区域，二是是否能够从损失函数层面进一步强化结构连续性约束。

首先考虑的是后处理。与前面几节的结构改进不同，后处理并不改变主干网络，也不直接影响特征提取过程，而是在模型输出的二值分割结果上执行规则修正。对于视网膜血管分割而言，较常见的问题之一就是出现面积较小、空间上相互孤立的连通域，它们往往来源于背景噪声、局部反光或误检。因此，本文采用基于面积阈值的小连通域筛除策略，希望在尽量保留主要血管结构的前提下，去除明显不合理的小区域。若将预测结果中的一个连通域记为 \(C_i\)，其像素面积记为 \(|C_i|\)，面积阈值记为 \(X\)，则后处理规则可以概括写为
\[
\hat{Y}_{\mathrm{post}}(C_i)=
\begin{cases}
1, & |C_i| \ge X,\\
0, & |C_i| < X.
\end{cases}
\]
其中，\(X=0\) 表示不做后处理；当 \(X\) 增大时，更多面积较小的连通域会被移除。图 \ref{fig:postprocess_threshold_cn_36} 展示了不同阈值下主要指标的变化趋势，表 \ref{tab:postprocess_cn_36} 给出了相应实验结果。

\begin{figure}[htbp]
\centering
\includegraphics[width=0.9\textwidth]{pic/Postprocess_Threshold_Metrics_dual_axis.png}
\caption{后处理面积阈值对主要指标的影响。左侧纵轴表示 SE、PR 和 F1-score，右侧纵轴表示 SP，用于展示不同量级指标随阈值变化的趋势。}
\label{fig:postprocess_threshold_cn_36}
\end{figure}

\begin{table}[htbp]
\centering
\caption{MS-DSConv + CA 在不同后处理阈值下的结果对比}
\label{tab:postprocess_cn_36}
\resizebox{\textwidth}{!}{
\begin{tabular}{lccccc}
\hline
实验组 & SE & SP & ACC & PR & F1-score \\
\hline
MS-DSConv + CA ($X=0$) & $0.8083$ & $0.9815$ & $0.9666$ & $0.7997$ & $0.8020$ \\
MS-DSConv + CA ($X=20$) & $0.8059$ & $0.9819$ & $0.9668$ & $0.8035$ & $0.8027$ \\
MS-DSConv + CA ($X=30$) & $0.8046$ & $0.9821$ & $0.9669$ & $0.8048$ & $0.8026$ \\
MS-DSConv + CA ($X=90$) & $0.7995$ & $0.9827$ & $0.9670$ & $0.8089$ & $0.8021$ \\
\hline
\end{tabular}
}
\end{table}

从表 \ref{tab:postprocess_cn_36} 可以看出，在不做后处理时，即 \(X=0\)，模型已经具有最强的血管检出能力，其中 SE 为 0.8083，F1-score 为 0.8020，但与此同时 PR 仅为 0.7997。随着阈值从 0 增加到 20、30 和 90，SP、PR 和 ACC 逐渐上升，而 SE 逐渐下降。这一变化提示，后处理能够清理掉一部分小噪声区域，但也可能删除部分真实的细小血管。进一步比较可以看到，在候选阈值中，\(X=20\) 时 F1-score 达到最高的 0.8027，因此本文将其保留为一种轻量修正，而不是高度调参后的后处理规则；当阈值继续增大时，结果虽然会更保守、更干净，但对细小血管的保留能力也随之下降。

在完成后处理实验之后，本文又进一步尝试从损失函数层面增强模型对血管拓扑连续性的约束。与输出端修正不同，clDice 关注的是训练阶段如何把骨架层面的连通性信息纳入优化目标 \cite{qi2023dsc}。在保持 MS-DSConv + CA 主体结构不变的前提下，本文在原有基础损失之上额外加入 clDice 损失，希望通过训练目标本身强化对细长结构骨架一致性和连通性的建模。若将预测结果记为 \(P\)，真实标注记为 \(G\)，则 clDice 可写为
\[
\mathrm{clDice}(P,G)=
\frac{2\,T_{\mathrm{prec}}(P,G)\,T_{\mathrm{sens}}(P,G)}
{T_{\mathrm{prec}}(P,G)+T_{\mathrm{sens}}(P,G)},
\]
其中，\(T_{\mathrm{prec}}\) 和 \(T_{\mathrm{sens}}\) 分别表示基于骨架重叠关系定义的拓扑精确率和拓扑灵敏度。根据当前实现，clDice 仅作用于最高分辨率输出分支，并以附加项的形式与原有损失共同优化。为便于与前一部分直接比较，本文同样在不同面积阈值 \(X\) 下评估加入 clDice 后的结果，具体见表 \ref{tab:cldice_cn_36}。

\begin{table}[htbp]
\centering
\caption{加入 clDice 后在不同后处理阈值下的结果对比}
\label{tab:cldice_cn_36}
\resizebox{\textwidth}{!}{
\begin{tabular}{lccccc}
\hline
实验组 & SE & SP & ACC & PR & F1-score \\
\hline
MS-DSConv + CA + clDice ($X=0$) & $0.8110 \pm 0.0477$ & $0.9806 \pm 0.0052$ & $0.9661 \pm 0.0041$ & $0.7944 \pm 0.0500$ & $0.8003 \pm 0.0243$ \\
MS-DSConv + CA + clDice ($X=20$) & $0.8079 \pm 0.0484$ & $0.9810 \pm 0.0052$ & $0.9663 \pm 0.0041$ & $0.7976 \pm 0.0500$ & $0.8004 \pm 0.0244$ \\
MS-DSConv + CA + clDice ($X=30$) & $0.8061 \pm 0.0486$ & $0.9812 \pm 0.0052$ & $0.9663 \pm 0.0041$ & $0.7990 \pm 0.0499$ & $0.8002 \pm 0.0245$ \\
MS-DSConv + CA + clDice ($X=90$) & $0.8002 \pm 0.0498$ & $0.9817 \pm 0.0051$ & $0.9663 \pm 0.0042$ & $0.8025 \pm 0.0496$ & $0.7989 \pm 0.0243$ \\
\hline
\end{tabular}
}
\end{table}

从表 \ref{tab:cldice_cn_36} 可以看出，加入 clDice 后，模型在不同阈值下的 SE 普遍略有提升。例如，在 \(X=0\) 时，SE 提升到 0.8110；在 \(X=20\) 时，SE 为 0.8079。这提示 clDice 作为面向细长结构拓扑一致性的附加损失，能够在一定程度上增强模型对细小血管和连续结构的敏感性。不过，与前一部分不加入 clDice 的结果相比，这一改动并未带来更好的综合性能。以当前表现最好的后处理阈值 \(X=20\) 为例，加入 clDice 后的 F1-score 为 0.8004，低于不加入 clDice 时的 0.8027；PR 也从 0.8035 下降到 0.7976。在当前实验条件下，clDice 更像是进一步放大了模型对潜在血管区域的敏感性，却没有同时把局部误检收住。

综合这一节可以看出，后处理和 clDice 分别代表了两种不同的补充思路：前者从输出端修正局部误检，后者从训练目标层面增强结构连续性约束。对本文当前阶段而言，限制最终结果的主要矛盾并不是“拓扑约束完全不够”，而是“高敏感结构模型带来的局部误检如何被有效修正”。因此，后处理在当前实验条件下带来了更直接且更稳定的收益，而 clDice 虽然是一项合理的探索，却未形成优于后处理方案的整体效果。基于这一判断，本文最终保留后处理作为后续结果优化的重要组成部分，而将 clDice 作为未保留的补充探索。

\subsection{本章小结}
本章围绕 nnU-Net 统一框架下的模型设计与实现过程，对从前期探索到后续结构优化的实验推进逻辑进行了梳理。前期实验表明，在 nnU-Net 既有预处理和训练流程基础上，额外输入增强与损失函数调整虽然会影响训练过程和局部指标表现，但并没有形成足够稳定的最终收益。尤其是在这一组前期实验设置下，输入端的额外预处理没有表现出预期中的明显增益。也正是在这一判断基础上，研究重心才逐步从训练层面的局部调整转向网络结构本身的改进。

后续研究先以 PA-Net 模块迁移作为第一轮结构探索的起点，通过完整迁移、单模块消融以及 AFFM 与深度监督关系的控制实验，对不同模块在当前框架中的适配性进行了初步筛选。结果显示，完整 PA-Net 迁移并不能形成更优方案，AFFM 的收益也并不稳定，而 LPT 则表现出更清晰的继续优化价值。围绕这一判断，后面的结构主线逐步收拢到 LPT，并沿着可变形采样、调制机制、多尺度聚合、位置注意力以及连续几何轨迹建模的方向继续推进。

随着结构连续性建模不断增强，模型在获得较高 SE 和 F1-score 的同时，也开始暴露出一定误检问题。因此，本章后半部分又从输出结果优化和损失函数约束两个层面进行了补充探索。比较结果显示，面积阈值后处理能够较直接地修正局部误检，并取得更稳定的综合收益；相比之下，clDice 虽然在方向上具有合理性，但在当前实验条件下并未带来更优的整体表现。对本章而言，最终被保留下来的并不是某一个孤立模块，而是一条经过逐步筛选后形成的主线方案及其必要的输出端补充策略。

本章主要回答了“模型如何逐步形成”这一问题。第二章侧重介绍相关方法基础，而第三章则进一步说明这些方法在本文中如何被取舍、组合和实现。在此基础上，下一章将进一步从整体结果、指标权衡和结构演化路径等角度，对本章形成的各阶段模型进行集中分析与讨论。

\section{实验结果与分析}

\subsection{整体结果对比}
为便于建立统一参照，本节选取本研究中具有代表性的阶段性模型进行整体比较。这里既包含直接迁移 PA-Net 的早期对照，也包含后续沿着 LPT 主线逐步演化得到的关键模型，用来勾勒本文性能变化的整体轨迹。讨论重点放在 SE、PR 与 F1-score 上，SP 和 ACC 则作为描述整体预测平衡的辅助指标。

\begin{table}[htbp]
\centering
\caption{第四章代表性模型整体结果对比}
\label{tab:chapter4_overall_cn}
\resizebox{\textwidth}{!}{
\begin{tabular}{lccccc}
\hline
模型 & SE & SP & ACC & PR & F1-score \\
\hline
nnU-Net v2 基线 & $0.7627$ & $0.9852$ & $0.9659$ & $0.8238$ & $0.7891$ \\
PA-Net (Full) & $0.7211$ & $0.9888$ & $0.9659$ & $0.8570$ & $0.7782$ \\
Deformable LPT (No AFFM, 4 point) & $0.7634$ & $0.9863$ & $0.9672$ & $0.8371$ & $0.7958$ \\
MS-Modulated (4 point) & $0.7582$ & $0.9870$ & $0.9673$ & $0.8419$ & $0.7954$ \\
MS-Modulated + CA & $0.7697$ & $0.9857$ & $0.9673$ & $0.8324$ & $0.7973$ \\
MS-DSConv + CA & $0.8083$ & $0.9815$ & $0.9666$ & $0.7997$ & $0.8020$ \\
MS-DSConv + CA ($X=20$) & $0.8059$ & $0.9819$ & $0.9668$ & $0.8035$ & $0.8027$ \\
\hline
\end{tabular}}
\end{table}

从表 \ref{tab:chapter4_overall_cn} 可以看到，直接迁移完整 PA-Net 并未形成优于 nnU-Net v2 基线的结果。虽然 PA-Net (Full) 的 PR 达到 0.8570，提示其输出相对保守、误检较少，但其 SE 仅为 0.7211，低于基线的 0.7627，最终 F1-score 也下降到 0.7782。完整迁移后的模型更倾向于保留高置信度血管区域，却牺牲了大量细小血管和弱响应区域的检出能力。这一现象提示，模块设计本身并非没有价值，但直接移植到 nnU-Net 框架中并不会自动带来更优结果。

与之相比，后续沿 LPT 路线展开的改进并不是让所有指标单调提高，而是更适合从 SE-PR 权衡的角度理解。Deformable LPT (No AFFM, 4 point) 将 F1-score 提升到 0.7958，提示在瓶颈层引入更灵活的采样方式对高层血管结构建模有帮助。随后，MS-Modulated 虽未继续提高 SE，但在 SP、ACC 与 PR 上更具优势，体现出更保守、更干净的预测风格；在其基础上加入 CA 后，SE 和 F1-score 再次上升，提示显式方向与位置信息仍然有助于细长血管恢复。进一步引入 MS-DSConv 后，模型的 SE 提升到 0.8083，F1-score 提升到 0.8020，提示连续几何轨迹建模比前一阶段离散采样形式更符合当前观察到的细小血管恢复趋势。

最终结果出现在 MS-DSConv + CA 配合面积阈值为 \(X=20\) 的后处理设置下，其 F1-score 达到 0.8027，相比基线的 0.7891 小幅提高 0.0136。这个增幅在绝对值上并不大，但它是在同一后期协议下得到的，并且与 SE-PR 权衡中的变化趋势相一致。该结果不是由后处理单独带来的，而是先依赖前面结构改进所建立的较强血管检出倾向，再通过轻量输出筛除机制修正低 PR 问题。最终保留方案更适合被理解为一系列结构选择和结果筛选共同形成的结果，而不是某一个孤立模块的单独作用。

\subsection{关键指标权衡分析}
从视网膜血管分割任务的需求出发，SE 主要反映细小血管和弱响应区域的检出能力，PR 更能体现模型对伪阳性的抑制能力，而 F1-score 则代表二者之间的综合平衡。因而，对各模型的判断不能只依据某一个指标是否更高，而应关注结构改动究竟把模型推向了更敏感还是更保守的方向。本节的图示主要聚焦于后期主线模型，即基线之后沿 LPT 路线逐步演化并最终保留的结构。

\begin{figure}[htbp]
\centering
\begin{minipage}{0.45\textwidth}
\centering
\includegraphics[width=\textwidth]{pic/comparion_of_se.png}\\
(a) SE 对比
\end{minipage}
\hfill
\begin{minipage}{0.45\textwidth}
\centering
\includegraphics[width=\textwidth]{pic/Comparison of F1-score.png}\\
(b) F1-score 对比
\end{minipage}
\caption{后期主线代表性模型在 SE 与 F1-score 上的变化趋势。图中依次展示了基线之后沿 LPT 主线演化得到的关键结构版本。}
\label{fig:chapter4_metric_pair_cn}
\end{figure}

\begin{figure}[htbp]
\centering
\includegraphics[width=0.82\textwidth]{pic/performace_balance_comparison.png}
\caption{后期主线代表性模型在灵敏度与精确率之间的性能平衡对比，其中包含最终后处理版本。}
\label{fig:chapter4_balance_cn}
\end{figure}

图 \ref{fig:chapter4_metric_pair_cn} 先给出了后期主线模型在 SE 和 F1-score 两项指标上的整体变化。可以看到，Deformable LPT 与 Modulated Deformable 都更偏向于增强细小血管检出。Deformable LPT 在基线基础上把 F1-score 提升到 0.7958，同时 SE 也略有提高，提示更灵活的采样方式可能有助于高层结构建模。Modulated Deformable 则进一步把 SE 提升到 0.7660。不过，从 F1-score 来看，这一步并没有比 Deformable LPT 带来更明显的综合优势，因此它更适合理解为沿着“增强检出能力”方向继续推进，而不是已经带来更好的整体平衡。

主线后半段的变化则更多体现在图 \ref{fig:chapter4_balance_cn} 所示的 SE-PR 关系上。MS-Modulated 的位置更靠右，对应较高 PR 和相对有限的 SE，这提示跨层多尺度采样更有利于压制背景噪声并维持较保守的预测。进一步加入 CA 后，模型的点位整体向上移动，表示在保持较高 PR 的同时，SE 得到回升。结合图 \ref{fig:chapter4_metric_pair_cn} 与表 \ref{tab:chapter4_overall_cn} 可以看到，这一步提高了 F1-score，同时 PR 仍保持在相对较高水平。CA 因此可以被理解为在已有多尺度特征交互基础上强化了与血管方向和位置信息相关的表达，而不只是单纯压低误检。

在此基础上，MS-DSConv + CA 体现出更鲜明的“以结构连续性换取更强检出能力”的倾向。图 \ref{fig:chapter4_metric_pair_cn} 显示，它取得了更高的 SE 和 F1-score，其中 SE 提升到 0.8083，F1-score 提升到 0.8020。图 \ref{fig:chapter4_balance_cn} 同时表明，这些变化伴随着 PR 下降。这提示 DSConv 可能使模型更加积极地识别潜在血管像素，这对于细小血管和连续结构恢复是有利的，但也可能带来更多小噪声区域。

图 \ref{fig:chapter4_balance_cn} 中最后一个点进一步体现了后处理的作用。对于 MS-DSConv + CA 而言，后处理并没有从根本上改变其“高敏感”的结构特征，而是通过去除部分孤立噪声，使 PR 从 0.7997 回升到 0.8035，同时只带来很小的 SE 损失。再结合表 \ref{tab:chapter4_overall_cn} 可以看到，后处理后的 F1-score 达到 0.8027，为当前阶段最高值。最终方案的优势并不在于某一项指标绝对最高，而在于它在血管检出能力、误检控制和整体平衡之间形成了较合适的折中。

\subsection{结构演化路径分析}

前两节主要讨论各模型的指标表现，本节进一步分析最终主线方案形成的原因。回到第三章可以看到，本文并不是先设定一个固定终点，再简单叠加不同模块，而是在多轮尝试之后逐步收窄改进方向。最早的 PA-Net 模块迁移已经提示，完整移植现有结构，并不会在当前框架下自然得到更优结果。真正显示出继续改进空间的，是位于瓶颈层位置的 LPT。后面的工作也因此逐步从解码端的融合问题，转回到血管结构本身：在高层特征里，怎样把弯曲、细长、分支频繁的几何特征表达得更合适。

沿着这一思路，Deformable LPT 和 Modulated Deformable 首先处理的是采样方式不够灵活的问题。固定的局部聚合在一般区域并不差，但放到走向变化快、尺度差异大的血管上，适应性仍然有限。引入可变形采样以后，特征聚合不再完全受固定位置约束；再加上调制权重后，模型会更倾向于把响应放在可能的血管区域上。对应到结果上，这一阶段最早显出来的变化主要是 SE 的上升。不过，这一步更多还是把“能不能更积极地把血管找出来”往前推进了一点，对浅层细节的补回还不够稳定，因此它还不能算是一个足够完整的答案。

再往后，从 Modulated Deformable 走到 MS-Modulated，再到 MS-Modulated + CA，关注点逐渐发生变化。问题不再只是采样是否足够灵活，而是来自不同层级的信息能否更有效地共同参与判断。多尺度采样的意义，也不只是多引入几层特征图参与计算，更关键的是把浅层细节重新带回高层结构判断里，让模型不至于只靠深层语义去推测血管走向。CA 的加入，则是在这个基础上继续补充方向和位置信息。走到这里，模型的改进已经不只是单纯追求更高的检出率，而是在尽量控制误检的前提下，使结构表达更加稳定。

\begin{figure}[htbp]
\centering
\includegraphics[width=\textwidth]{pic/overlay_case_comparison_24_34.png}
\caption{代表性病例的分割结果可视化对比。自左至右分别为原图、人工标注、nnU-Net 基线、MS-Modulated + CA、MS-DSConv + CA 及其后处理结果。}
\label{fig:chapter4_overlay_cn}
\end{figure}

图 \ref{fig:chapter4_overlay_cn} 给出的两组代表性病例，更适合用来观察不同阶段结果的风格差异，而不是直接判断中间几种结构的绝对优劣。整体上看，基线结果相对保守，后续模型对潜在血管区域的响应更加积极。Case 24 中，MS-Modulated + CA 相比基线并没有表现出更明显的细分支恢复，整体更接近一种较紧的预测；换成 MS-DSConv + CA 之后，模型保留了更多血管区域，因此与标注的重合度更高。Case 34 则显示出另一点：中间模型之间并不是线性地此消彼长。相较于基线，MS-Modulated + CA 和 MS-DSConv + CA 都有改进，但侧重点不完全一样。前者看起来更干净一些，后者对部分血管区域保留得更积极，最后和标注的重合度也略高。至于最后一列的后处理，其作用更接近结果清理，把零散的小区域收掉，并没有改写前一列已经形成的主干趋势。

因此，DSConv 最后被保留下来，并不只是因为它把某几个指标继续往上推了一点。更重要的是，它和前面几步改动之间存在较清晰的衔接关系：多尺度采样和 CA 先提供跨层特征与位置信息基础，DSConv 再把这种偏向细长连续结构的建模方式落实到局部采样层面。这样一来，模型面对弯曲明显、分支较多的血管时，会更愿意给出连续响应。其代价也比较清楚，即可能带来更多零散噪声，这也是后处理仍然具有保留价值的原因。从整条演化路径看，最后落到 MS-DSConv + CA，再配合后处理，并不是因为它堆叠的模块最多，而是因为这一组合更能回应当前任务中的主要问题。

\subsection{后处理与 clDice 的补充分析}

上一节已经说明，主线最后落在 MS-DSConv + CA，并不是因为走到这里以后问题已经全部解决，而是因为在这一阶段，它在结构表达和细小血管检出上的收益最值得保留。与此同时，模型的短板也变得更清楚了。结合图 \ref{fig:chapter4_balance_cn} 和表 \ref{tab:chapter4_overall_cn} 可以看到，MS-DSConv + CA 的 SE 和 F1-score 都已经比较高，但 PR 仍然偏低。此时影响最终结果的主要问题，已经不再是主干血管能不能分出来，也不单纯是拓扑约束够不够，而是模型在更积极恢复血管区域的同时，也带出了一些零散噪声和局部伪阳性。

在这种情况下，后处理仍然是最终方案的一部分。它并不替代前面的结构改进，只是在输出端对已经暴露出来的误差做一次轻量修正。对 MS-DSConv + CA 这类偏高敏感的模型来说，问题通常不在主干走向完全错误，而在于结果边缘残留了一些面积较小、彼此孤立的连通域。这类误差更适合直接在输出结果上清理，而不是再回到主干网络里进一步压制。第三章表 \ref{tab:postprocess_cn_36} 也支持这一判断：在 \(X=0\) 时，模型的 F1-score 为 0.8020，PR 为 0.7997；加入 \(X=20\) 的面积阈值后，PR 回升到 0.8035，F1-score 进一步提高到 0.8027，而 SE 只从 0.8083 小幅下降到 0.8059。这个提升幅度不大，但正好落在当前最需要修正的误差类型上，因此本文最终版本保留后处理。

clDice 的情况就不太一样了。它没有被保留下来，不是因为完全没有作用，而是因为它改善的方向和当前阶段最主要的问题没有完全对上。按照第三章表 \ref{tab:cldice_cn_36} 的结果，加入 clDice 之后，模型在一些设置下的 SE 略有上升。例如，在 \(X=0\) 时，SE 从 0.8083 提高到 0.8110。这提示，把骨架连通性信息放进训练目标里，会让模型更愿意对细长结构作出响应。可问题也在这里：这种变化并没有转化为更好的整体平衡。以更有代表性的 \(X=20\) 设置为例，加入 clDice 后，F1-score 为 0.8004，低于未加入时的 0.8027；PR 也从 0.8035 下降到 0.7976。在当前实验条件下，clDice 更像是进一步放大了模型对潜在血管区域的敏感性，却没有同时把局部误检收住。

综合来看，后处理和 clDice 对应的是两种不太一样的补充思路。前者发生在输出端，处理的是已经形成的局部错误；后者发生在训练阶段，试图从优化目标本身去推动结构连续性。对本文这条主线来说，到了 MS-DSConv + CA 之后，结构连续性本身已经不是最薄弱的一环，更突出的反而是高敏感预测之后的结果清理问题。在这个前提下，后处理的收益会更直接，也更稳定；clDice 虽然方向合理，但没有形成优于后处理方案的综合表现。因此，本文最终保留下来的，是“MS-DSConv + CA + 轻量后处理”这一组合，而 clDice 则更适合被看作一项有启发意义、但没有进入最终方案的补充尝试。

\subsection{本章小结}

本章将前文形成的各阶段模型放到同一参照下重新比较，讨论的不只是结果高低，也包括这些结果是怎样逐步出现的。整体比较可以看出，完整迁移 PA-Net 并没有在当前框架下超过基线，持续显示出优化价值的主要是沿 LPT 展开的主线。随着可变形采样、调制机制、多尺度建模、位置注意力以及连续几何轨迹建模依次加入，模型对细小血管和结构连续性的表达逐渐增强，MS-DSConv + CA 也因此成为本章分析中最值得保留的后期结构。

不过，第四章的重点并不是证明“哪一个模块单独最好”。把几组指标和可视化结果放在一起看，后续结构改动一直围绕同一组权衡展开：在检出能力、误检控制和整体平衡之间反复调整。前面的 Deformable LPT 和 Modulated Deformable，更偏向把血管尽量找出来；再往后的多尺度和 CA，让结果慢慢稳下来；到了 MS-DSConv + CA，细长连续结构的恢复又被进一步加强，同时也把局部误检的问题一并带了出来。所以最后留下来的，不是某个孤立模块，而是一条经过多轮比较后逐渐收拢的主线方案。

后处理和 clDice 的补充分析，则进一步强化了这一判断。对当前这套模型来说，更需要解决的问题已经不是“能不能分出血管”，而是“分出来以后如何清理多余的小噪声”。面积阈值后处理正好对上了这个问题，因此能提供更直接、也更稳定的收益；clDice 虽然方向合理，也在一定程度上增强了模型对细长结构的敏感性，但没有换来更好的综合平衡。归结起来，第四章最后保留下来的，是以 MS-DSConv + CA 为核心、再辅以轻量后处理的结果优化方案。

\section{总结与展望}

\subsection{全文总结}
本文围绕基于 nnU-Net 的视网膜血管分割方法展开研究，重点关注细小血管恢复、结构连续性和误检控制。本文没有脱离现有成熟框架重新搭建一套全新模型，而是选择以 nnU-Net 作为统一基础，在相对稳定的实验条件下逐步考察不同结构改动和结果优化策略到底能带来什么。

全文实验分为两个阶段。前期探索显示，额外输入增强和损失函数调整会影响收敛过程，但在该设置下没有带来明确的最终收益。后期主实验因此转向统一五折协议下的结构改进。经过 PA-Net 迁移、模块筛选、可变形与多尺度采样、Coordinate Attention 以及 Dynamic Snake Convolution，后期主线逐渐收敛到 MS-DSConv + CA。

最终保留版本由 MS-DSConv + CA 和轻量连通域后处理组成。在后期统一协议下，其 F1-score 达到 0.8027，比后期 nnU-Net 基线的 0.7891 高 0.0136。这个提升在绝对值上并不大，但与观察到的 SE-PR 权衡关系一致：结构改动与更好的血管恢复相关，而后处理修正了一部分残留的局部伪阳性。

\subsection{主要工作与结论}
本文的主要工作和相应结论，大致可以归纳为以下几个方面。

第一，前期探索显示，在本文设置下，常规输入增强和损失调整不足以在 nnU-Net 基础上形成稳定最终收益。这也推动了后续研究重心转向结构设计。

第二，完整迁移 PA-Net 并没有优于后期 nnU-Net 基线。这一轮筛选更有价值的地方在于，它把后续方向收敛到瓶颈层位置的 LPT，而不是直接沿用整套外部结构。

第三，后期 LPT 路线主要通过改变血管恢复与误检控制之间的平衡来影响性能。可变形采样和 DSConv 使模型对细长结构更敏感，而多尺度建模、CA 和轻量后处理则帮助输出结果避免过度噪声化。

第四，最终保留结果是 MS-DSConv + CA 配合 \(X=20\) 后处理。它的优势应理解为统一协议下的小幅稳定改进，而不是一次很大的性能跃升。clDice 具有参考价值，但没有被保留；它在部分设置下提高了细长结构敏感性，但在当前实验中相较后处理方案降低了 PR 和 F1-score。

\subsection{不足与改进方向}
尽管本文在 nnU-Net 框架下形成了一条相对清晰的后期主线方案，但仍有一些限制需要说明。

首先，本文对 nnU-Net 内部预处理流程的任务适配分析仍不够充分。前期实验表明，在图像进入网络前额外加入输入端预处理，并没有带来明显收益，这在一定程度上说明 nnU-Net 自身已经具备较为完善的内部预处理机制。但与此同时，本文并没有继续拆解其内部预处理各步骤在当前任务中的具体作用，也没有围绕视网膜血管分割的特点，对归一化方式、通道组织形式以及低对比度区域增强等环节作更有针对性的调整。因此，当前对“为什么额外预处理没有产生明显效果”的理解，仍主要停留在实验现象层面，还缺少更细化的流程分析与机制验证。

其次，本文当前的验证范围仍然有限。虽然前期探索和后期主实验分别建立了相对清晰的比较参照，后期实验也在统一协议下完成了较为系统的结构比较，但整体结果仍主要围绕当前采用的公开数据条件展开。这样的安排有助于保证横向比较的可解释性，却也意味着本文对模型泛化能力的讨论仍然比较有限。尤其是面对不同成像设备、不同病理表现以及不同标注习惯下的数据时，当前保留下来的方案是否仍然能够保持同样的优势，还需要在更丰富的数据条件下进一步验证。

再次，本文最终方案仍然依赖轻量后处理来修正局部伪阳性区域。这说明主线结构虽然与更好的细小血管恢复和连续结构表达相关，但在误检控制方面，仍然没有完全通过端到端方式解决问题。当前结果中的最终平衡，一部分来自结构建模本身，另一部分则依赖输出端的额外筛除。这样的组合在当前设置下有效，但从方法完整性来看，如何把这类结果修正能力进一步吸收到网络结构或训练目标内部，仍然是一个值得继续推进的问题。

最后，本文对模型有效性的解释仍以实验结果分析为主。当前结论主要建立在整体指标变化、局部可视化结果以及不同结构之间的对比关系之上，能够说明“哪些改动有用、哪些改动未被保留”，但对其中一些现象的内部机制解释还不够深入。例如，不同结构为什么会在某些病例上呈现出更敏感或更保守的预测风格，DSConv 与多尺度采样、位置注意力之间究竟是如何相互作用的，这些问题目前仍缺少更细的中间特征分析和可解释性支撑。此外，从定性结果来看，当前保留下来的模型虽然对主要血管结构和细小血管检出有一定改善，但与人工标注相比，仍可见部分细小分支遗漏、局部断裂以及弱响应区域恢复不足等现象。这也说明，本文方法在更细粒度的结构完整性恢复上仍有继续提升的空间。也正因如此，本文的结论更适合被看作当前实验条件下的阶段性结果，而不是对相关机制的完全封闭解释。

\subsection{后续研究展望}
结合本文当前工作基础，后续研究可以沿以下几个方向继续展开。

一条比较直接的后续方向，是围绕 nnU-Net 内部预处理流程做更细的任务适配分析。前期实验表明，额外加入输入端预处理并没有带来明显收益，但这并不意味着相关环节已经没有继续研究的空间。更值得继续做的，是把 nnU-Net 内部预处理各步骤在视网膜血管分割任务中的具体作用拆开来看，再结合归一化方式、通道组织形式以及低对比度区域处理等因素，探索更贴合当前任务特点的预处理改进方案。

另一个值得继续推进的问题，是更端到端的误检控制方式。本文最终保留下来的方案通过轻量后处理较好地改善了局部伪阳性问题，但这类修正毕竟仍然发生在输出端。如果后续能够把类似的结果清理能力进一步融入网络结构或训练目标内部，就有可能在保持细小血管恢复能力的同时，减少对额外规则筛除的依赖，让整套方法更完整一些。

结构设计本身也还有进一步打磨的空间。本文当前保留下来的主线结构虽然在细小血管恢复和连续结构表达方面取得了较好结果，但多尺度采样与相关结构增强主要还是集中在瓶颈层及编码端特征聚合一侧，并没有在解码端形成与之对应的镜像式设计。从这个角度看，当前主线虽然仍保留了 U-Net 的基本编码器—解码器框架，但在具体改动分布上已经表现出一定的不对称性。后续研究可以继续往前走一步，看看在保持现有主线有效性的基础上，是否能够构建出更协调的对称式改进结构，并进一步比较对称化设计对模型性能、结构表达能力以及结果稳定性的实际影响。这个方向的意义，不只是为了让结构形式上更完整，也有助于重新检验当前主线中哪些改动是真正关键的。

除此之外，还需要在更丰富的数据条件下继续验证当前方案的适用范围。本文目前的主要结论仍建立在现有公开数据和统一实验协议基础之上，后续如果能够进一步扩展到不同成像设备、不同病理表现以及不同标注风格的数据条件中，会更有助于判断模型的泛化能力和实际应用潜力。

\clearpage
\bibliographystyle{IEEEtran}
\bibliography{references}

\end{document}

